%-------------------------------%
%                               %
%    Hauptdokument              %
%    Diplomarbeit               %
%    Nicole Kikillus            %
%                               %
%    November 2003 - April 2004 %
%                               %
%-------------------------------%

\documentclass{wissdoc} %\textheight 23.0cm \textwidth 15.5cm
%\topmargin0cm \evensidemargin0cm \oddsidemargin0cm
\setlength{\parindent}{0pt}
\usepackage{a4,ngerman}
%\usepackage[a4paper, left=2.7cm, right=2.3cm, top=3.8cm, bottom=2.6cm]{geometry}
%\usepackage[a4paper, left=2.8cm, right=1.7cm, top=3.5cm, bottom=2.4cm]{geometry}
% ORG zum 11.08.05 \usepackage[a4paper, left=2.4cm, right=2.0cm, top=3.7cm, bottom=2.6cm]{geometry}
% ORG zum 18.07.05 \usepackage[a4paper, left=2.4cm, right=2cm, top=3.8cm, bottom=2.6cm]{geometry}
% ORG zum 06.07.05 \usepackage[a4paper, left=2.5cm, right=2.1cm, top=3.8cm, bottom=2.6cm]{geometry}
% ORG \usepackage[a4paper, left=3.6cm, right=3.2cm, top=3.8cm, bottom=2.6cm]{geometry}
%\usepackage[a4paper, left=2.5cm, right=2.5cm, top=1.2cm, bottom=2.3cm]{geometry}
\usepackage{graphicx}
\usepackage{epsfig}
\usepackage{float}
\usepackage[utf8]{inputenc}
\usepackage{url}
\usepackage{makeidx}
%\usepackage{natbib}
% \usepackage[dvips]{color}
\usepackage{color}
\usepackage{fancyhdr}
\usepackage{color}
\usepackage{pstricks}
\usepackage{rotating}
\usepackage[footnotesize,bf]{caption2}
\usepackage{eurosym}
\usepackage{pdfpages}
\usepackage{amsmath}
\usepackage{listings}
 \definecolor{middlegray}{rgb}{0.5,0.5,0.5}
 \definecolor{lightgray}{rgb}{0.8,0.8,0.8}
 \definecolor{orange}{rgb}{0.8,0.3,0.3}
 \definecolor{yac}{rgb}{0.6,0.6,0.1}
\lstset{numbers=left, numberstyle=\tiny, numbersep=5pt}
\lstset{language=python}
 \lstset{
   basicstyle=\small\ttfamily,
   keywordstyle=\bfseries\ttfamily\color{orange},
   stringstyle=\color{blue}\ttfamily,
   commentstyle=\color{middlegray}\ttfamily,
   emph={square}, 
   emphstyle=\color{blue}\texttt,
   emph={[2]root,base},
   emphstyle={[2]\color{yac}\texttt},
   showstringspaces=false,
   flexiblecolumns=false,
   tabsize=2,
   numbers=left,
   numberstyle=\tiny,
   numberblanklines=false,
   stepnumber=1,
   numbersep=10pt,
   xleftmargin=15pt
 }
\makeatletter \renewcommand{\fps@figure}{htbp} \makeatother
\renewcommand{\textfraction}{.1} 
\setlength{\parindent}{0em}
\renewcommand{\baselinestretch}{1.2}
%\KOMAoptions{twoside,BCOR=5mm,DIV=11}
\KOMAoptions{twoside,BCOR=5mm,DIV=12,footinclude=false,headinclude=true}

\recalctypearea

% Graphiken sollen nicht so oft auf Leeren Seiten stehen
\renewcommand{\topfraction}{0.9}
\renewcommand{\textfraction}{0.05}

% Macros, nicht unbedingt notwendig
\input{macros}

\newcommand{\blankpage}{% Leerseite ohne Seitennummer, nächste Seite rechts
 \clearpage{\pagestyle{empty}\cleardoublepage}
}


\pagestyle{headings}

%% *********************************
%% Formel- und Abkürzungsverzeichnis
%% *********************************

%% Achtung: an Kommandozeile folgender Befehl ausführen:
%%          makeindex -o diplom.gls -s nomencl.ist diplom.glo

\usepackage{nomencl}
\usepackage{ifthen}
\renewcommand{\nomname}{\chapter*{Glossar}}
\makeglossary
\makenomenclature
\setlength{\nomlabelwidth}{.25\hsize}
%\renewcommand{\nomlabel}[1]{\emph{#1}}
\setlength{\nomitemsep}{-\parsep}
%\makeglossary
%\makenomenclature


%% Quellcodeverzeichnis
\renewcommand\lstlistlistingname{Quellcodeverzeichnis}

%% Auskommentieren einer dieser Zeilen bewirkt, daß nur diese
%% eine Datei übersetzt wird

%\includeonly{Titel}
%\includeonly{Dank}
%\includeonly{Zitat}
%\includeonly{Kap_1}
%\includeonly{Kap_2}
%\includeonly{Kap_3}
%\includeonly{Kap_4}
%\includeonly{Kap_5}
%\includeonly{Kap_6}
%\includeonly{Anh_a}
%\includeonly{Anh_b}
%\includeonly{Anh_c}
%\includeonly{Anh_d}
%\includeonly{Anh_e}
%\includeonly{Anh_f}
%\includeonly{Anh_g}
%\includeonly{Tempverlauf}

%% **************
%% Dokumentbeginn
%% **************

%%% Überschrift 4ter Ordnung
\setcounter{secnumdepth}{4}

\usepackage[noindentafter]{titlesec}

% paragraphs sehen aus wie subsubsubsections
\titleformat{\paragraph}[hang]{\bf}{\thetitle\quad}{0pt}{}                      
\titlespacing{\paragraph}{0pt}{1em}{0.5em} 

% subparagraphs sehen aus wie vorher paragraphs
\titleformat{\subparagraph}[runin]{\bf}{}{0.5em}{}
\titlespacing{\subparagraph}{0pt}{1em}{1em}


%%%-----------------------------------
\begin{document} 
%%%-------------------------
\begin{center}
%\subsubsection*{}%\vspace{3cm}
{\epsfig{file=Grafik/unilogo.eps, scale=0.25}}
\end{center}
\begin{center}
\textsf{\large UNIVERSIT�T KARLSRUHE (TH) - FORSCHUNGSUNIVERSIT�T}\\
			\vspace{0.5cm}
\textsf{\large FAKULT�T F�R ELEKTRO- UND INFORMATIONSTECHNIK}\\
			\vspace{0.5cm}
\textsf{\large INSTITUT F�R BIOMEDIZINISCHE TECHNIK}\\
			\vspace{2cm}
\textsf{{\Huge DIPLOMARBEIT\\}}
			\vspace{3cm}
\textsf{\Large Neuartige Verfahren der bewegungslosen\vspace{0.25cm}\\\Large Respirationsdetektion anhand des Elektrokardigramm}\\
			\vspace{3.5cm}
\begin{table}[ht]
	\begin{tabular}{p{0cm} p{5.1cm} l}
		& \textsf{\large vorgelegt von} 								& \textsf{\large cand. el. Philipp Scheidel}\vspace{0.5cm}\\

		& \textsf{\large Betreuer}  & \textsf{\large Prof. Dr. Armin Bolz}\vspace{0.5cm}\\
		
		& \textsf{\large }			& \textsf{\large Dipl.-Ing. Marc J�ger}\vspace{0.5cm}\\
		
		& \textsf{\large Abgabetermin}				& \textsf{\large 31.09.2005}\\
	\end{tabular}	
\end{table}
\end{center}
\thispagestyle{empty}
\clearpage
%\begin{center}
    \begin{quote}
    Tell me and I'll forget. \\
    Show me and I may remember.\\
    Involve me and I'll understand.
    \end{quote}
    Old Chinese Proverb
\end{center}



%\thispagestyle{empty}
\chapter*{} 
\pagenumbering{Roman}% Seitennummerierung durch arabische Ziffern
\vspace*{10.5cm} {\setlength{\parindent}{0pt}\textbf{Eidesstattliche Erklärung}\\
Hiermit erkläre ich an Eides Statt, dass ich die vorliegende Diplomarbeit selbständig und ohne unzulässige fremde Hilfsmittel angefertigt habe. Die verwendeten Literaturquellen sind im Literaturverzeichnis vollständig angegeben. Die Arbeit wurde in gleicher oder ähnlicher Form keiner anderen Prüfungsbehörde zur Erlangung eines akademischen Grades vorgelegt.\\ 

Karlsruhe, 31.10.2010\\

\chapter*{Abstract}
\begin{small}
In dieser Arbeit wird ein kabelloses Sensornetzwerk zur Überwachung von Patienten bei einem 
Massenanfall von Verletzten entworfen und implementiert. Als Grundlage dieser Entwicklung
dient der am IBT entwickelte \emph{Erste-Hilfe-Sensor}, der mit einem neuartigen kapazitiven
Verfahren Puls und Atmung eines Patienten überwacht.
Hierzu wird zunächst eine Untersuchung aller in Frage kommenden Funktechnologien vorgenommen,
wobei ein besonderes Augenmerk auf Energieffizienz und Stabilität der Funktechnologie gelegt wird.
Hierbei stellt sich \emph{ZigBee} als am besten geeignetes Protokoll heraus, und es wird entschieden,
eine Umsetzung auf Basis eines \emph{ZigBee}-Moduls des Herstellers \emph{Atmel} durchzuführen. 
Dieses Modul wird einer genaueren Betrachtung unterzogen, und eine Reihe von Schwierigkeiten werden
identifiziert. Für diese Schwierigkeiten werden allgemeingültige Lösungen entwickelt. Auf Basis dieser
allgemeinen Lösungen werden später spezielle Lösungen für die \emph{Java-Plattform} sowie für den
Mikrocontroller des \emph{Erste-Hilfe-Sensors} entwickelt.
\\                    
Bei der Entwicklung des Gesamtsystems wird ein komponentenbasierter Ansatz gewählt. Hierdurch wird nicht
nur die Austauschbarkeit der einzelnen Komponenten erreicht, sondern auch deren Verteilung über
mehrere Rechner ermöglicht. Für die Verbindung der einzelnen Komponenten zueinander werden 
\emph{Corba}-Schnittstellen verwendet. Diese Schnittstellen sind so gestaltet, dass auch externe
Software angebunden werden kann. Die Anbindung exemplarisch am Beispiel einer \emph{Weboberfläche} 
durchgeführt.
Außerdem wird eine Platine entwickelt, auf der das \emph{ZigBee}-Modul und der
\emph{Mikrocontroller} des \emph{Erste-Hilfe-Sensors} integriert sind. Diese Platine wird zur 
Entwicklung eines \emph{ZigBee}-Treibers für die Firmware des \emph{Erste-Hilfe-Sensors} verwendet.
\\
Das Gesamtsystem wird zum Abschluß eingehend evaluiert. Hierbei wird gezeigt, dass alle Komponenten
erfolgreich miteinander interagieren. Außerdem wird die Leistungsaufnahme des \emph{ZigBee}-Moduls
untersucht, und die Abhängigkeit von den Parametern des Energiesparmodus berechnet. Hierbei ergibt 
sich ein Strom von ca. 1,86~mA bei der Wahl von realistischen Parametern für den 
Energiesparmodus. Dies entspricht einer Leistungsaufnahme von ca 5,6~mW.
%\begin{small} -- Falls es nicht auf eine Seite passt darf die Schrift im Abstract kleiner gemacht werden, aber nur beim Abstract!
%... Die Zusammenfassung ("`Abstract"') ist nach dem Titel der zweitwichtigste Bestandteil einer
%wissenschaftlichen Arbeit. Sie sollten deshalb für die Zusammenfassung, ebenso wie für den
%Titel, besonders viel Mühe und Zeit verwenden, da die gesamte wissenschaftliche Arbeit nur
%von sehr wenigen Wissenschaftlern gelesen wird, die Zusammenfassung aber von vielen. Aus
%der Zusammenfassung muß hervorgehen, wovon die Arbeit handelt, worauf sie aufbaut, und
%vor allen Dingen, welche neuen Erkenntnisse gewonnen wurden. Die wichtigsten Ergebnisse
%der Arbeit müssen kurz und präzise aufgezählt werden. Es genügt nicht zu schreiben, daß dies
%und jenes in der Arbeit behandelt werden. Wichtig ist, daß die "`harten Fakten"', welche sich
%aus den Untersuchungen ergeben haben, aufgelistet sind.
%Handelt es sich um eine theoretische Arbeit, dann müssen Sie erwähnen, von welchen
%Gleichungen Sie ausgegangen sind und welche Näherungen Sie verwendet haben; bei einer
%experimentellen Arbeit müssen Sie erwähnen, welche Experimente Sie durchgeführt haben
%und eventuell auch, welche Auswertemethoden (falls nicht Standardmethoden) Sie verwendet
%haben.Beachten Sie, daß Ihre Arbeit von Wissenschaftlern unterschiedlicher Herkunft und
%Ausbildung gelesen wird. Bedenken Sie, daß sich auch Wissenschaftler für Ihre Arbeit
%interessieren können, die aus benachbarten Disziplinen stammen und nicht mit dem von Ihnen
%verwendeten wissenschaftlichen "`Jargon"' vertraut sind, oder solche, welche die in Ihrem Fach
%üblichen Abkürzungen nicht kennen. Deshalb soll die Zusammenfassung für alle
%(natur-)wissenschaftlich gebildeten Leser verständlich sein. Das bedingt, daß eventuell
%benutzte Abkürzungen erklärt werden müssen, und daß nur solche Begriffe vorkommen
%dürfen, die ein "`normaler"' Wissenschaftler üblicherweise kennt oder die er notfalls in einem
%Lexikon nachschlagen kann. 
%Die Zusammenfassung sollte keine Literaturhinweise enthalten. Die Zusammenfassung ist ein
%selbständiger Teil der Arbeit. Das bedeutet, daß die in der Zusammenfassung erklärten
%Abkürzungen im Hauptteil noch einmal erklärt werden müssen.
%Einerseits darf die Zusammenfassung nicht zu lang sein (max eine Seite), andererseits muß sie aber auch alle wichtigen
%Informationen über Ihre Untersuchungen enthalten. Auf präzise Formulierungen ist größten
%Wert zu legen. 
\end{small}

%%%------------------------------------------------VORWORT----------------------------------------------------------------------------
%\thispagestyle{empty}
\chapter*{Vorwort} 

Diese Arbeit entstand in Rahmen des \emph{Erste-Hilfe-Sensor}-Projektes in Kooperation mit Herrn Jan Tepelmann,
der zur selben Zeit eine Studienarbeit geschrieben hat\cite{Jan}. Die Arbeit von Herrn Tepelmann
beschäftigt sich mit dem Entwurf und der Implementierung der Überwachungssoftware des Sensornetzwerkes
(\emph{MANVSuite}), wohingegen sich die hier vorliegende Arbeit mit dem Entwurf der Hardware und der Interaktion der
einzelnen Komponenten des Sensornetzwerkes befasst.\\
\\
Der \emph{Erste-Hilfe-Sensor} verwendet ein neuartiges Verfahren zur Überwachung von Patienten mit Hilfe eines induktiven
Verfahrens. Dies bietet einige Vorteil gegenüber klassischen \emph{EKG} basierten Verfahren. Die eigentliche Funktionsweise
des \emph{Erste-Hilfe-Sensor} ist für die vorliegende Arbeit nicht weiter relevant, lediglich auf den verwendeten
Mikroprozessor muss Rücksicht genommen wurden.\\
\\
Bei dieser Arbeit wurde ich von zahlreichen Helfern unterstüzt, bei denen ich mich auf diesem Weg bedanken möchte. 
Die wichtigste Person bei der Durchführung dieser Arbeit war mein Projektpartner Herr cand. inform Jan Tepelmann, der 
auf meine Bitten bereit war, die Implementierung der Überwachungssoftware in Form einer Studienarbeit durchzuführen. Er
war immer ein wertvoller Diskussionpartner, ohne dessen Hilfe es in der gegebenen Zeit nicht möglich gewesen, 
diese Arbeit in der vorliegenden Form anzufertigen.\\
Mein besonderer Dank gilt meinen beiden Gutachtern und meinem Betreuer: Herr Professor Dr. Armin Bolz hat durch seine 
spannende Vorlesung das Interesse an der biomedizinischen Technik erst bei mir geweckt, und sich sehr viel Zeit 
genommen, mich zu dieser interdisziplinären Arbeit zu ermutigen. Herr Professor Dr. Rüdiger Dillmann hat sich bereit 
erklärt, von Seiten der Fakultät für Informatik ein Gutachten durchzuführen -- ohne dies wäre ein Durchführen 
dieser Arbeit garnicht möglich gewesen. Mein Betreuer, Herr Dr. Marc Jäger hat mir immer wieder mit wertvollen Gesprächen
zur Seite gestanden, und mir viele Tipps zur korrekten Durchführung gegeben. Insbesondere hat er mir als externem 
Diplomand ein hohes Maß entgegengebracht, und mir so ermöglicht, das geplante Praktisch umzusetzen. Seine Firma,
die Neocor GmbH hat dabei die Kosten für die notwendige Hardware übernommen.\\
Von unschätzbarem Wert während des Studiums war und ist mein langjähriger guter Freund Herr Juniorprofessor Dr Christoph
Sorge, der sich nicht nur die Zeit zum Korrekturlesen genommen, sondern mir bereits seit vielen Jahren wertvolle
Tipps und Hilfestellungen ums Studieren gegeben hat, und sich immer wieder Zeit für meine Sorgen und Nöte nimmt. 
Eine weitere sehr wertvolle Hilfe war Herr cand. inform. Alexander Neumann, der vor vielen Jahren die Lust auf 
Mikrocontroller bei mir geweckt hat, und mir während der Diplomarbeit ein wichtiger Diskussionspartner für 
Hardwarefragen war. Auch meinem Freund Herrn Dipl.-Inform. (FH) Markus Müller möchte ich für das Korrekturlesen danken.
Sehr herzlich wurde ich von den Mitarbeitern des IBT aufgenommen. Hierbei bedanke ich mich insbesondere bei Herrn
M.-Eng Daniel Wettach, der mir bei Fragen zum \emph{ADuC-Microcontroller} immer weitergeholfen hat, sowie Herrn
Dipl.-Ing Stefan Fernsner für die Hilfe beim Erstellen der \emph{MANVNode}-Platinen. Auch bedanken möchte ich mich bei 
Frau Dipl.Phys. Kerstin Grimmel, Herrn Dipl.-Ing Nikolas 
Lentz und Herrn M. Sc. Firas Salih für zahlreiche anregende und hilfreiche Diskussionen.\\
Von ganzem Herzen bedanken möchte ich mich bei meiner Lebenspartnerin Frau Dipl.-Phys. Jennifer Girrbach sowie meinen Eltern,
meinen Geschwistern und meiner Familie und Freunden für den Rückhalt den sie mir in dieser anstregenden Zeit gegeben, und den
Mut, den sie mir zugesprochen haben.

%%%%------------------------------------------------DANKSAGUNG-------------------------------------------------------------------------
\chapter*{Danksagung}
An dieser Stelle m�chte ich den Personen danken, ohne die diese Arbeit so nicht entstanden und auch nicht m�glich gewesen w�re. Ein herzliches Dankesch�n gilt...
\begin{itemize}
\item ... 
\item ... 
\item ...  
\item ... 
\item ... 
\item ... 
\end{itemize}

Die Danksagung ist ein wichtiger Teil der Arbeit. Hier sollten Sie sich
bei all denen bedanken, die Ihnen bei den beschriebenen Forschungen behilflich waren. Es ist
zwingend erforderlich, da� in der Danksagung steht, welche Institutionen Ihnen evtl. Daten
kostenlos zur Verf�gung gestellt haben und von welchen Institutionen Ihre Arbeit finanziert
wurde. Allgemein ist zu sagen, da�
man in der Danksagung nicht geizen sollte. Durch das Erw�hnen von Hilfen, die man erhalten
hat, kann man sich T�ren �ffnen, die sp�ter eventuell von gro�em Nutzen sein k�nnen. Ein
weiterer Tip: Schicken Sie jedem/jeder, der/die in der Danksagung erw�hnt wird, ein
Exemplar Ihrer Arbeit (am besten mit Widmung). Er/sie wird sich freuen.


%%%------------------------------------------------INHALTS-,ABBILDUNGS-,TABELLENVERZEICHNIS-------------------------------------------
\setcounter{tocdepth}{4}
\tableofcontents 
\listoffigures
\addcontentsline{toc}{chapter}{Abbildungsverzeichnis}
\listoftables
\addcontentsline{toc}{chapter}{Tabellenverzeichnis}
\lstlistoflistings
\addcontentsline{toc}{chapter}{Quellcodeverzeichnis}

%%%------------------------------------------------GLOSSAR-----------------------------------------------------------------------------

%\chapter{Glossar}
\markboth{Glossar}{}
\printnomenclature
\addcontentsline{toc}{chapter}{Glossar}




%%%------------------------------------------------EINLEITUNG-------------------------------------------------------------------------
\chapter{Einleitung}
\pagenumbering{arabic} % Seitennummerierung durch arabische Ziffern
\section{Motivation der Arbeit}
...

\section{Aufgaben und Ziele der Arbeit}
...
\section{Gliederung und Vorgehensweise der Arbeit}
...

%%%------------------------------------------------Grundlagen------------------------------------------------------------------
\chapter{Grundlagen}\label{Grundlagen}

\section{Scatternet}

\section{Kabellose Übertragungsprotokolle}
    \subsection{Einführung}
    In diesem Abschnitt werden die gängisten Funkprotokolle kurz vorgestellt. Insbesondere wird erläutert,
    inwieweit das entsprechende Protokoll als Grundlage für das zu entwickelnde Sensornetz geeignet ist.

    Mit Ausnahme von DECT und GSM bzw. UMTS ist diesen Protokollen gemein, dass sie sich alle im ISM-Band 
    befinden. 

\subsection{DECT}

Bei DECT ("`Digital Enhanced Cordless Telecommunications"') handelt es sich um einen Standard, 
der vor allem zur Anbindung von Schnurlostelefonen an eine Basisstation gedacht ist\footnote{Es gibt
jedoch auch weitere Anwendungen wie z.B. Babyfone.}. 

In Europa wird einn eigenes Frequenzband im Bereich von 1800 bis 1900 MHz verwendet, in dem 10 
Kanäle zur Verfügung stehen. Pro Kanal können maximal 32kbit Nutzdaten pro Sekunde übertragen
werden. Die maximal zulässige Sendeleistung beträgt 250mW, womit eine Reichweite von ca. 30-50
Metern in Gebäuden und ca. 300m im Freien realisiert werden kann. Jede Basisstation kann
bis zu 6 Geräte anbinden.

Beim Einsatz außerhalb Europas muss bedacht werden, dass die Verwendung der Frequenzen von
1800 bis 1900 MHZ hier evtl. nicht zulässig ist. In diesem Fall muss auf das ISM-Band 
ausgewichen werden, welches sich hier mit anderen Anwendungen geteilt werden mus.

DECT bietet eine optionale Verschlüsselung der Nutzdaten, welche jedoch im Jahr 2009 geknackt
wurde, so dass DECT mittlerweile als unsicher gelten muss.

Aufgrund der geringen Nutzdatenmenge sowie der Einschränkung auf 6 Teilnehmer ist
DECT für den Einsatz als Sensornetz-Protokoll nicht geeignet.

\subsection{GSM/UMTS}


\subsection{WLAN}\label{wlan}
\subsection{Übersicht}
WLAN oder Wi-Fi bezeichnet den heute gängigen Standard eines Funkprotokolls zum Aufbau
von kabellosen lokalen Netzwerken. Es gibt mehrere Versionen des Standards, die verbreitetsten
sind IEEE 802.11a, IEEE 802.11b/g und IEEE 802.11n. 

Es sind zwei Betriebsmodi möglich:

\begin{itemize}
    \item{Infrastruktur-Modus:} Eine zentrale Station ("`Access-Point"') dient als Basisstation
                                für alle weitere Stationen. Jede Station, die am Netzwerk
                                teilnimmt, muss hierzu die Signale des \textsl{Access-Points} 
                                empfangen können.
    \item{Ad-hoc-Modus:} Diese Betriebsmodus kommt ohne zentrale Komponente aus. Es wird eine
                         Peer-to-Peer-Verbindung zwischen allen am Netzwerk teilnehmenden 
                         Stationen aufgebaut. Hierzu ist es notwendig, dass alle Stationen
                         sich gegenseitig empfangen können.
\end{itemize}

Der Infrastruktur-Modus bietet gegenüber des Ad-Hoc-Modus' klare Vorteile: Im Gegensatz
zum Ad-hoc-Modus ist es nicht notwendig, dass alle Stationen sich gegenseitig empfangen
müssen, es reicht aus, wenn der Acces-Point empfangen werden kann. Hierdurch ist eine
größere Ausbreitung des Netzwerks möglich als im Ad-Hoc-Modus. 

Es werden verschiedene Datenübertragungsraten unterstützt. Standardmässig wird
immer die größtmögliche Übertragungsrate gewählt, die Störungsfrei verwendet
werden kann. Je weiter sich die Stationen voneinander entfernen, desto geringer
wird die Übertragungsrate, bis schliesslich die niedrigst mögliche Übertragungsrate von
1MBit/sec erreicht wird.

Es gibt eine ganze Menge von Unterstandards, die sich durch zulässige Frequenzen, 
Übertragungsraten und Sendeleistung unterscheiden. Teilweise ist auch Kanalbündelung 
vorgesehen. Nicht jeder Standard kann in jedem Land eingesetzt werden, und die meisten
Endgeräte unterstützen nur eine Teilmenge dieser Stanrdards. Hier sollen nur die
wichtigsten drei Standards erwähnt werden.

\subsubsection{IEEE 802.11b}
Bei IEEE 802.11b handelt es sich um den ältesten WLAN-Standard, der bereits 1999 
spezifiziert wurde. Dieser Standard wird praktisch von jedem WLAN fähigen 
Endgerät unterstützt. Die Kommunikation findet im ISM-Band im Bereich von 2.4GHz
statt. Je nach Land sind 11-14 Kanäle möglich, die sich jedoch teilweise überlappen.
Dies führt dazu, dass maximal 3 Netzwerke ohne Störungen gleichzeitig betrieben werden
können. Die maximale Sendeleistung beträgt 100mW. Es sind Übertragungsraten von 
5,5 bis 11 MBit (brutto) möglich. Die Nettoübertragungsrate beträgt ca. 50\% der
Bruttorate. Es sind Reichweiten bis 40m (Innen) bzw. 100m (Außen) möglich\cite{WirelessNetworking}.

\subsubsection{IEEE 802.11g}
Der IEEE 802.11g Standard stellt eine Erweiterung des IEEE 802.11b Standards da.
Wesentliche Neuerung ist eine Erhöhung der Bruttodatenrate von 11 auf 54MBit/sec,
von denen netto ca. 40\% zur Verfügung stehen. Erwähnenswert ist, dass die
beiden Standards Interoperabel sind, d.h. ein 802.11b Gerät kann einem
802.11g Netzwerk beitreten und umgekehrt. Dies ist auch der Grund, weshalb
dieser Standard momentan am weitesten verbreitet ist.

\subsubsection{IEEE 802.11a}
Der IEEE 802.11a-Standard verwendet Frequenzen im 5GHz Bereich. Er ist daher
inkompatibel zum IEEE-802.11b/g-Standard. Je nach Frequenzband sind 
Sendeleistungen zwischen 30 und 1000mW zulässig. Mit dem passenden Frequenzband
sind daher höhere Reichweiten als mit dem IEEE-802.11b/g-Standard möglich.
Die Bruttodatenrate beträgt bis zu 54MBit/sec. 
Ein Vorteil des IEEE-802.11a Standards ist die Kommunikation im 5GHz Bereich.
Aktuell ist dieser Bereich noch wenig genutzt, so dass in diesem
Bereich oft ein störungsärmerer Betrieb als im 2,4GHz-Bereich möglich ist.
Es ist jedoch zu erwarten, dass sich dies in Zukunft ändern wird.

\subsubsection{Leistungsaufnahme}
Der Energiebedarf für WLAN ist relativ hoch. Beispielsweise benötigt der
vom Hersteller Broadcom als besonders energiesparend bezeichnete Chip
BCM4326 bis zu 100mA zum Empfangen und zwischen 141 und 190mA zum Senden.
\footnote{BCM4326 Datasheet}. 

\subsubsection{Anzahl Teilnehmer}

\subsection{WPAN: Wireless Personal Area Networks}
    \subsubsection{Übersicht}
        Als WPANs ("`Wireless Personal Area Networks"') werden kabellose Kleinnetzwerke bezeichnet, die dazu dienen,
        wenige Geräte über kurze Entfernungen (mehrere Meter) miteinander zu verbinden. Sie dienen als Ersatz 
        von Kabelverbindungen zur Anbindung von Peripherie an Computergeräte (z.B. zur Verbindung von
        Headsets mit Mobiltelefonen oder von Tastatur und Maus mit einem PC).
        

    \subsubsection{IEEE 802.15}
        Der IEEE 802.15-Standard behandelt \textsl{Wireless Personal Area Networks}. Er ist in mehrere Unterstandards
        aufgeteilt:

         \begin{itemize}
            \item{IEEE 802.15.1:} Bluetooth 1.2
            \item{IEEE 802.15.2:} Zusammenarbeit zwischen IEEE 802.15 (WPAN) und IEEE 802.11 (WLAN)
            \item{IEEE 802.15.3:} WPANs mit hohen Datenübertragungsraten (20MBit/sec und höher)
            \item{IEEE 802.15.4:} WPANs mit niedriger Datenübertragungsraten
        \end{itemize}

        Für diese Arbeit sind vor allem der Bluetooth und der ZigBee Standard interessant. 

    \subsubsection{IEEE 802.15.1: Bluetooth}
        \paragraph{Überblick}
            Bluetooth wurde ursprünglich von dem Mobilfunkhersteller Ericcson als Ersatz für
            RS-232 Verbindungen entwickelt. Die Entwicklung des Bluetooth Standards erfolgt
            heute unter der Regie der \textsl{Bluetooth Special Interest Group} ("`SIG"').
            Version 1.1 des Bluetooth Standard wurde 
            von der IEEE als IEEE 802.15.1-2002 übernommen. Nach Veröffentlichen einer
            weiteren Version IEEE 802.15.1-2005, die dem Bluetooth 1.2 Standard entspricht,
            wurde von der IEEE jedoch beschlossen, nicht weiter mit der \textsl{Bluetooth SIG}
            zu kooperiern, so dass es keine weiteren Versionen des IEEE 802.15.1 Standards 
            geben wird. Aktuell ist Version 4.0 des Bluetooth Standard, wobei jedoch die
            meisten Geräte nur geringere Standards (Typischerweise 2.0 oder 2.1) unterstützen.
            Für den Bluetooth 4.0 Standard existiert zum Zeitpunkt dieser Arbeit keine 
            Implementierung auf dem Markt.\\
            Die wichtigsten Meilensteine der Bluetooth-Entwicklung kann folgendermaßen
            zusammengefasst werden:

            \begin{itemize}
                \item{Bluetooth 1.1:} Erste Version von praktischer Relevanz. Entspricht 
                                      IEEE 802.15.1-2002.
                \item{Bluetooth 1.2:} Entspricht IEEE 802.15.2-2005. Bringt einige Verbesserungen
                                      gegenüber der Version 1.1 wie z.B. schnelleres Finden
                                      von Endgeräten (Discovery), höhere Störsicherheit durch
                                      die Verwendung von AFH\footnote{Frequency-hopping spread spectrum},
                                      Übertragungsraten bis 721kbit/sec.
                \item{Bluetooth 2.0:} Einführung des EDR\footnote{Enhanced Data Rate}-Modus mit bis zu
                                      3.0 MBit/sec (2.1 MBit/sec netto).
                \item{Bluetooth 2.1:} Vereinfachung des Pairings (vgl. \ref{BlueetoothPairing}) durch 
                                      Einführung von SSP\footnote{Secure Simple Pairing}, Verbesserung
                                      der Sicherheit durch explizite Aushandlung der Verschlüsselung.
                \item{Bluetooth 3.0:} Einführung eines Hochgeschwindigkeits-Datenkanals auf Basis von
                                      IEEE 802.11 (vgl. \ref{wlan}) mit bis zu 24MBit/sec., verbessertes
                                      Powermanagement, Einführung von Verbindungslosen Datentelegrammen
                                      (Unicasts).
                \item{Bluetooth 4.0:} Einführung des \textsl{Blueetooth Low Energy} Standards (vgl. \ref{wibree}).
            \end{itemize}

        \paragraph{Pairing}\label{BluetoothPairing}
        \paragraph{Reichweite}
        \paragraph{Übertragungsrate}
        \paragraph{Leistungsaufnahme}
        \paragraph{Anzahl Teilnehmer}
    \subsubsection{IEEE-802.15.4: ZigBee}
    \subsubsection{Wibree: Bluetooth Low Energy}\label{wibree}
    \subsection{Weitere Protokolle}
        - WiMAX
        - Mikrowellen-Richtfunk
        - Z-Wave
        - Wireless USB

    \subsection{Diskussion}

\section{Der Analog Devices ADuC702X Mikrocontroller}
\section{Java}
\section{Corba}


%%%------------------------------------------------STAND DER TECHNIK------------------------------------------------------------------
\chapter{Stand der Technik}\label{Stand der Technik}

\section{Allgemeine kabellose Technologien zur Patientenüberwachung}
Auf dem Markt ist eine breite Palette an kabellosen Produkten zur Patientenüberwachung vorhanden. Die meisten davon setzen
entweder \emph{Bluetooth} oder proprietäre Protokolle wie \emph{Z-Wave} zur Funkübertragung ein, \emph{ZigBee} ist bisher in
keinem Medizinprodukt auf dem Markt zu finden. Fast alle dieser Systeme sind dazu konzipiert, lediglich einen einzelnen 
Patienten zu überwachen. Die Funktechnologie dient hier vor allem dazu, die Anzahl der Kabel zwischen Patient und Monitoren 
zu reduzieren, um so die Mobilität des Patienten zu erhöhen. Besonders häufig sind kabellose EKG-Geräte, Pulsoximeter oder 
invasive Blutdruckmessgeräte anzutreffen. 
-- FIXME: Edit here --

\section{Spezielle Technologien zur Patientenüberwachung im MANV-Szenario}

\chapter{Analyse}
    \section{Anforderungen an ein kabelloses System zur Überwachung von Patienten}
    Austauschbarkeit beachten!

    \section{Kommunikation mit dem \emph{ZigBit}-Modul}\label{analyse_kommunikation_mit_zigbit_modul}
        \subsection{Problemstellung}
            Das \emph{SerialNet} Protokoll ist grundsätzlich synchron. Ein Programm, dass mit
            einem \emph{ZigBit}-Modul kommuniziert lässt sich daher am einfachsten nach folgendem
            Muster aufbauen:

            \begin{lstlisting}{frame=analyse1}
def sendeBefehl(befehl)
    # Senden des Befehls
    serialPort.send(befehl)

    # Antwort vom ZigBit-Modul lesen.
    # (Blockiert so lange, bis Antwort empfangen wurde)
    statusString = serialPort.readLine()

    # Antwort interpretieren
    if statusString == 'OK':
        return True
    
    else:
        return False
            \end{lstlisting}
           
        
            Bei dieser Vorgehensweise ergeben sich jedoch zwei Probleme:

            \begin{enumerate}
                \item{Asynchron auftretende Ereignisse:} Die Abarbeitung eines \emph{SerialNet}-Befehls kann 
                          bis zu einer halben Minute dauern\footnote{Dies tritt insbesondere dann auf, wenn eine
                          synchrone Datenübertragung zu einem entfernten Knoten durchgeführt werden soll, der sich
                          selbst gerade im Energiesparmodus befindet. Der sendende Knoten ist hierbei so lange
                          blockiert, bis entweder der Empfang der Daten quittiert wurde, oder aber ein 
                          Timeout aufgetreten ist}. Während das Steuerungprogramm auf die Quittierung des 
                          Befehls mit einer Status-Meldung wartet, können asynchrone Ereignisse auftreten.
                          Im oben abgebildeten Programm würde dieses Ereigniss fälschlicherweise als Antwort
                          auf den gesendeten Befehl interpretiert werden. Dies hätte zur Folge, dass nun 
                          Steurungsprogramm und \emph{ZigBit}-Modul nicht mehr synchron zueinander ("`out-of-sync"') 
                          sind.

                \item{Komplexe Rückgabewerte:} Bestimmte Befehle der \emph{SerialNet}-Firmware liefern nicht nur
                         einen einfachen Statuscode sondern eine komplexe, Mehrzeilige Antwort zurück. Auch hier ergibt
                         sich bei obigem Programm wieder das Problem, dass dieses bereits nach der ersten Zeile
                         fälschlicherweise davon ausgehen würde, dass die Antwort abgeschlossen wäre, obwohl noch weitere
                         Zeilen folgen können. Auch in diesem Falle wären Steurungsprogramm und \emph{ZigBit}-Modul nichtmehr
                         synchron zueinander.
            \end{enumerate}

        \subsection{Behandlung asynchron auftretender Ereignisse}
            \label{behandlung_asynchron_auftretender_ereignisse}

            Um asynchrone Ereignisse behandeln zu können, müssen diese zunächst erkannt werden. Dies kann mit
            Hilfe einer Liste aller Ereignisse und Statuscodes geschehen. In diesem Falle wird so lange
            von dem \emph{ZigBit}-Modul gelesen, bis ein Statuscode empfangen wurde. Zwischendurch empfangene
            Ereignisse können entweder direkt behandelt werden, oder zunächst in einer Liste zwischengespeichert
            werden, um sie zu einem späteren Zeitpunkt zu behandeln. Beide Verfahren haben sowohl Vor- als auch
            Nachteile:

            \begin{itemize}
                \item{Direktes Behandeln von Events:} Dieses Verfahren hat den Vorteil, dass auf empfangene 
                Ereignisse sofort reagiert werden kann; so lässt sich z.B. ein empfangener Alarm direkt
                auf dem Bildschirm darstellen und es muss nicht erst auf die Beendigung des vorherigen
                Befehls gewartet werden. Erfordert das behandelte Event den Zugriff auf das 
                \emph{ZigBit}-Modul um z.B. den Empfang einer Nachricht zu quittieren, so ist eine
                besondere Vorsicht erforderlich, da sonst entweder Verklemmungen ("'Deadlock"`) oder
                ein Verlust der Synchronisierung auftreten kann\footnote{Eine mögliche Lösung ist die
                Verwendung einer Warteschlange, in die die zu versendenden Daten eingereiht werden.}.
                Dieses Verfahren eignet sich besonders für Systeme, in denen auf eingehende Ereignisse besonders
                schnell reagiert werden muss.

                \item{Einreihen von Ereignisse in eine Warteschlange:} Diese Verfahren hat den Vorteil, dass
                es sehr einfach umgesetzt werden kann. Alle Ereignisse, die während der Bearbeitung eines Befehls
                auftreten werden zunächst in eine Warteschlange eingereiht. Mit der Abarbeitung der Ereignisse
                wird so lange gewartet, bis der laufende Befehl abgeschlossen wurde. Während der Abarbeitung der
                Ereignisse muss dann keine Rücksicht mehr darauf genommen werden, in welchem Zustand sich das
                \emph{ZigBit}-Modul befindet, da es sich nach Abarbeitung des Befehls sicher im Betriebsbereiten
                Zustand befindet\footnote{Ausnahme hiervon ist, wenn zwischenzeitlich die Netzwerkverbindung verloren
                gegangen ist, was natürlich gesondert behandelt werden muss.}. Nachteilig ist, dass Ereignisse
                erst behandelt werden, wenn Befehle komplett abgearbeitet wurden. Bei Ereignissen, zu deren Bearbeitung
                nicht auf das \emph{ZigBit}-Modul zugegriffen werden muss bedeutet dies eine unnötige Verzögerung.
                Dieses Verfahren eignet sich daher eher für Systeme, in denen deutlich mehr Befehle als Ereignisse
                auftreten.
            \end{itemize}

            Durch die Verwendung von Multithreading kann darüber hinaus eine Kombination aus beiden Verfahren
            verwenden werden. In diesem Fall wird von einem Thread aus die Bearbeitung der Befehle, inkl. des
            Befüllens der Ereigniswareteschlange vorgenommen, während von einem anderen Thread aus diese Warteschlange
            parallel dazu abgearbeitet wird.\\
            \\
            Das folgende Pseudocode-Programm zeigt exemplarisch das Behandeln von Events über das Einreihen in
            Warteschlangen:

            \begin{lstlisting}{frame=analyse2}
eventQueue = []

def sendeBefehl(befehl)
    # Senden des Befehls
    serialPort.send(befehl)

    # Antwort vom ZigBit-Modul lesen.
    # (Blockiert so lange, bis Antwort empfangen wurde)
    statusString = serialPort.readLine()

    # Events behandeln 
    while statusString not in ('OK', 'ERROR'):
        eventQueue.append(statusString)
        statusTring = serialPort.readLine()

    # Status Code interpretieren
    if statusString == 'OK':
        return True
    
    else:
        return False
            \end{lstlisting}

        \subsection{Behandlung von komplexen Antworten}
            \label{behandlung_von_komplexen_antworten}
            Bei der Behandlung von komplexen Antworten kommt der Umstand zur Hilfe, dass bereits
            im Vorraus bekannt ist, welche Befehle eine komplexe Antwort erzeugen. Bei welchen Befehlen dies
            der Fall ist kann aus Tabelle~\ref{zigbit_befehle} entnommen werden. Im Steuerprogramm kann beim
            Parsen des Befehls daher entsprechend Rücksicht darauf genommen werden, ob eine einzeilige oder
            eine mehrzeilige Antwort erwartet wird. Komplexe Antworten werden immer mit einem Statuscode
            beendet, d.h. wenn eine komplexe Antwort erwartet wird, muss so lange vom \emph{ZigBit}-Modul
            gelesen werden, biss ein Statuscode empfangen wird. Auch hierbei muss natürlich darauf geachtet
            werden, dass zwischen einzelnen Zeilen einer Antwort Ereignisse auftreten können, die entsprechend
            der in Abschnitt~\ref{behandlung_asynchron_auftretender_ereignisse} vorgestellten Lösung behandelt 
            werden müssen.

            \begin{lstlisting}{frame=analyse3}
def sendeBefehl(befehl)
    ## Senden des Befehls
    serialPort.send(befehl.getCmd())

    if befehl.resultIsComplex():
        ## Komplexe Antwort wird erwartet
        ## So lange lesen, bis ein Status-Code empfangen wurde
        complexResult = []
        statusString = ''

        while statusString not in ('OK', 'ERROR'):
            statusString = serialPort.readLine()
            complexResult.append(statusString)

        return statusString

    else:
        ## Antwort vom ZigBit-Modul lesen.
        ## (Blockiert so lange, bis Antwort empfangen wurde)
        statusString = serialPort.readLine()

        ## Antwort interpretieren
        if statusString == 'OK':
            return True
        
        else:
            return False

            \end{lstlisting}

        \subsection{Kombination beider Lösungen}

        Kombiniert man die Lösungen aus Abschnitt~\ref{behandlung_asynchron_auftretender_ereignisse} und 
        Abschnitt~\ref{behandlung_von_komplexen_antworten} ergibt sich ein weiteres Problem: Bei der
        Lösung aus Abschnitt~\ref{behandlung_asynchron_auftretender_ereignisse} wird davon ausgegangen, 
        dass alle Ergebnisse, bei denen es sich nicht um Statuscodes handelt, Ereignisse sind. Diese Annahme
        ist jedoch hinfällig, sobald komplexe Antworten vom ZigBit Modul empfangen werden. Daher muss hierbei
        jede Antwort einer genaueren Prüfung unternommen werden, ob es sich hierbei um einen Status Code, 
        ein Event oder eine Komplexe Antwort handelt. Hierbei kann eine einfache Mustererkennung verwendet
        werden.


            \begin{lstlisting}{frame=analyse4}
eventQueue = []

def getResultType(input):
    if input in ('OK', 'ERROR'):
        return "Status code"
    elif input[0:4]  == 'DATA' or \
         input[0:11] == 'EVENT:JOINED' or\ 
         input[0:10] == 'EVENT:LOST':
        return "Event"
    else:
        return "Complex result"

def sendeBefehl(befehl)
    ## Senden des Befehls
    serialPort.send(befehl.getCmd())

    if befehl.resultIsComplex():
        ## Es wird eine komplexe Antwort erwartet

        ## Von seriellem Port legen
        statusString = serialPort.readLine()

        ## Komplexe Antwort wird erwartet
        ## So lange lesen, bis ein Status-Code empfangen wurde
        complexResult = []
        
        while getResultType(statusString) != 'Status code':
            if getResultType(statusString) == 'Event':
                eventQueue.append(statusString)
            else:
                complexResult.append(statusString)

            statusString = serialPort.readLine()

        ## Ende Erreicht. Nun noch den statusCode anhaengen und
        ## Ergebnis zurueckliefern
        complexResult.append(statusString)
        return complexResult

    else:
        ## Antwort vom ZigBit-Modul lesen.
        ## (Blockiert so lange, bis Antwort empfangen wurde)
        statusString = serialPort.readLine()

        # Events behandeln 
        while statusString not in ('OK', 'ERROR'):
            eventQueue.append(statusString)
            statusTring = serialPort.readLine()

        # Status Code interpretieren
        if statusString == 'OK':
            return True
        
        else:
            return False

            \end{lstlisting}


    \section{Powermanagement}\label{analyse_powermanagement}
        Die Spannungsversorgung des Erste-Hilfe-Sensors erfolgt über Batterie. Da der Wechsel der Batterie
        während eines MANVs umständlich und zeitraubend ist, ist eine möglichst hohe Batterielaufzeit
        wünschenswert. Am besten wäre hierbei, wenn überhaupt kein Batteriewechsel notwendig wäre. Im
        empfansbereiten Zustand braucht das \emph{ZigBit}-Modul laut Datenblatt einen Strom von 19mA.
        Zur Verringerung der Leistungsaufnahme bietet das \emph{ZigBit}-Modul einen Energiesparmodus,
        indem lediglich der interne Speicher gehalten und alle andere Hardware abgeschaltet wird. In
        diesem Modus braucht das Modul laut Datenblatt weniger als $6\mu{}A$.\\
        \\
        Zum Stromsparen wird das Modul periodisch für eine bestimmte Zeit in den
        Energiesparmodus versetzt. Während dieser Zeit kann das Modul weder von der Steurungssoftware angesprochen
        werden, noch können Daten vom Netzwerk empfangen werden. Damit in dieser Zeit keine Informationen
        verloren gehen, werden vom nächstgelegenen Router alle Nachrichten an dieses Modul zwischengespeichert,
        bis dieses wieder im Empfangsbereiten Zustand ist. Bei Verwendung der \emph{SerialNet}-Firmware wird
        die Dauer des Verweilens im Energiesparmodus über den Befehl \emph{AT+WPWR} konfiguriert. Die Konfiguration
        erfolgt in Schritten zu 100ms.\\
        \\
        Zum Aufruf des Energiesparmodus bieten sich nun zwei Möglichkeiten:

        \begin{itemize}
            \item{Automatsisches Aufrufen des Energiesparmodus:} Hierbei wird über den zweiten Paramter
            des Befehls \emph{AT+WPWR} ein Intervall angegeben, in dem der Energiesparmodus periodisch
            aufgerufen wird. Hierbei erfolgt die Eingabe in Schritten von 10ms.

            \item{Manuelles Aufrufen des Energiesparmodus:} Alternativ zum automatischen Aufrufen kann
            der Energiesparmodus über den Befehl \emph{AT+WSLEEP} auch manuell aufgerufen werden. Dies
            kann entweder entweder Alternativ zum automatischen Aufrufen oder als Ergänzung verwendet
            werden.

        \end{itemize}

        Bei der Verwendung des Energiesparmoduses muss sichergestellt werden, dass keine Informationen,
        die über die \emph{UART}-Schnittstelle an das Modul gesendet werden, verloren gehen. Hierzu 
        bietet die \emph{SerialNet}-Firmware die Möglichkeit, über die \emph{CTS}-Leitung der 
        \emph{UART}-Schnittstelle ihre Empfangsbereitschaft zu signalisieren. Die Steurungssoftware
        auf der Gegenseite muss hierzu vor jedem Senden den Zustand der \emph{CTS}-Leitung
        kontrollieren, und darf nur senden, wenn der Zustand den Wert 0 hat. Damit die
        \emph{SerialNet}-Firmware diesen Zustand korrekt kommuniziert, muss über den Befehl \emph{AT+IFC}
        die Flußkontrolle aktiviert werden. \\
        \\
        Zu beachten ist, dass die Konfiguration des Energiesparmodueses bei allen Teilnehmern des Netzes
        (also auch auf Routern und Coordinator) gleich sein sollte. Der eigentliche Energeisparmodus steht
        allerdings nur auf Endknoten, also \emph{RFD}s zur Verfügung.



%%%------------------------------------------------Entwurf -------------------------------------------------------------
\chapter{Entwurf}

\section{Gesamtsystem}
\Abbildungps{!htb}{1}{diagramme/Gesamtuebersicht.pdf}{label}{Überblick über das Gesamtsystem}

\section{Hardware}
\subsection{MANVNode}
\subsubsection{Beschreibung}
\Abbildungps{!htb}{1.25}{diagramme/platine_manvnode-crop.pdf}{label}{Entwurf der Platine der MANVNode}
Die MANVNode ist eine Test- und Entwicklungsplatine zum Test der Sensornetzes. Auf ihr befindet sich ein 
ADuC-7026-Mikrocontroller der Firma Analog Devices sowie ein Atmel ZigBit-Funkmodul des Typs ATZB-24-A2. Die Platine 
dient zur Simulation des Erste-Hilfe-Sensors. An Stelle einer echten Messung generiert die Platine Messwerte mit Hilfe 
eines Zufallszahlen-Generators.
Dier Generator kann in drei Zuständen beschrieben werden:

\begin{itemize}
    \item{Grün:} Simulation eines Patienten mit gutem Kreislaufzustand. Puls zwischen 60 und 120 Schlägen pro Minute, 
                 Atemzüge pro Minute zwischen 12 und 15. Periodischer Versand der Meldung "`Zustand OK"'.
    \item{Gelb:} Simulation eines Patienten mit beginnend kritischem aber noch nicht akut Lebensbedrohlichem 
                 Kreislaufzustand. Puls zwischen 120 und 180 Schlägen pro Minute, Atmung zwischen 60 und 80
                 Zügen pro Minute (Hyperventilation).
    \item{Rot:}  Simulation eines Patienten in kritischem Zustand mit Puls zwischen 0 und 10 (entsprechend
                 Kammerflimmern oder Kreislaufstillstand) sowie Atmung zwischen (0 und 2) mit bestehndem
                 akutem Handlungsbedarf. Priorisierte Generierung eienr Alarmmeldung.
\end{itemize}

Zur Einstellung und Wechsel der Zustände befinden sich 3 Schalter an der Platine. Der Zustand wird mit Hilfe von
drei farbigen LEDs visualisiert. Zusätzlich kann über einen Piezzo-Summer ein Alarm auch akustisch wiedergegeben werden.
Durch die Kombination der drei LEDs werden zusätzlich Status- und Fehlerinformationen der Platine während des 
Initialisierungsvorgangs ausgegeben.

\subsection{Spannungsversorgung}
Um die Platine möglichst vielseitig einzusetzen wurde wert auf eine möglichst flexible Spannungsversorgung gelegt. Diese
soll möglichst Robust gegenüber Spannungsschwankungen sein. Kern der Spannungsversorgung ist ein 
Low-Dropout-Spannungsregler des Typs LP2980-3.3 der Firma National Semiconductor. Dieser liefert eine stabile
3,3V Spannung bei einem maximalen Strom von 50mA. Die Eingangsspannung kann zwischen 2.1V bis 16V variiert werden.
Die Spannungsversorgung kann damit z.B. über USB (5V-Versorgungsspannung), ein externes Netzteil oder eine 
9V-Blockbatterie gespeist werden. Durch die hohe maximale Versorgungsspannung von 9V sind Schäden durch versehentliches
anschließen eines falschen Netzteils weitestgehend ausgeschlossen. Die Spannungsversorgung ist zusätzlich mit einem 10mF 
Kondensator stabilisiert.

Zusätzlich wurde die die Spannungsversorgung weiter stabilisiert. Hierzu wurde jeder Spannungseingang aller Bauteile der
Platine mit einem 100nF Kondensator gepuffert. Diese sind möglichst nahe an den entsprechenden Versorgungspins angebracht.
Kurzzeitige Spannungsabfälle, wie sie z.B. beim Schalten von LEDs, des Piezzo-Summers oder des Aufwachsens des ZigBee-Moduls
aus dem Stromsparmodus auftreten können, werden so zuverlässig abgefangen.

\subsection{Programmierschnittstellen}

\subsubsection{Serielle Schnittstellen}
Die \emph{MANVNode}-Platine verfügt über mehrere seriellen Schnittstellen. Im äußeren Bereich der Platine ist eine serielle 
Schnittstelle (JP1) vorhanden, welche der Standardbelegung des Programmierkabels des ADuC-7026-Evaluationsboards besitzt.
Bereits vorhandene Programmierkabel können somit ohne weiteres zur Programmierung des Mikrocontrollers verwendet werden.
Zusätzlich besthet die Möglichkeit, auf die serielle Schnittstelle des ZigBit-Moduls zuzugreifen, worüber das Herunterladen 
von Firmware auf das Modul, sowie der Einstellung bestimmter Parameter wie der Netzwerkkennung des Gerätes möglich ist.
Der Zugriff erfolgt ebenfalls über die Stiftleiste JP1, allerdings muss hierfür RX- und TX-Leitung des Programmierkabels 
verdreht, und alle Jumper des Schalters S1 unterbrochen werden. 

\subsubsection{JTAG}
Die \emph{MANVNode}-Platine verfügt neben seriellen auch über eine JTAG-Schnittstelle. Diese ermöglich einen Zugriff mit
gängigen Programmierwerkzeugen wie z.B. Eclipse, Crossworks oder Kyle $\mu{}$Vision. Über die Schnittstelle kann sowohl 
Programmierung als auch Debugging des ADuC-7026-Mikrocontrollers durchgeführt werden. Ein Zugriff auf das ZigBit
Modul ist allerdings nur indirekt möglich\footnote{Einige Werkzeuge bieten die Möglichkeit, eine Consolenverbindung
über den UART des ADuC-7026-Mikrocontrollers zu öffnen, so direkt mit dem ZigBit-Modul kommunikation möglich ist.}.
Ein Zugriff über JTAG ist nicht vorgesehen.

\subsection{LEDs}

\subsubsection{Beschreibung der einzelnen LEDs}
Die \emph{MANVNode}-Platine verfügt insgesamt über 6 LEDs, die zur Anzeige des aktuellen Zustands sowie zur Fehlerdiagnose dienen.
\begin{itemize}
    \item{PWR:} Diese LED befindet sich direkt in der Spannungsversorgung der Platine. Diese LED leuchtet, sobald
                die Platine mit Spannung versorgt wird. 
    \item{GREEN, YELLOW, RED}: Diese LEDs sind direkt an den Mikrocontroller angeschlossen und visualisieren den Zustand des
                               Zufallszahlengenrators entweder durch dauerhaftes Leuchten (GREEN, YELLOW) oder im Alarmzustand
                               durch schnelles Blinken (RED). 
    \item{BLUE}: Diese LED signalisiert, ob eine Alarmunterdrückung ("`Mute"') besteht.
    \item{ZIGBIT-STATUS}: Diese LED leuchtet, sobald das ZigBit-Modul über das CTS-Signal ("`Clear to send"') die Bereitschaft, 
                          Daten zu empfangen, signalisiert.
\end{itemize}
                 
    

\subsubsection{Diagnosecodes}
Neben den oben dargestellten einfachen Signalisierungen hat die Firmware der \emph{MANVNode}-Platine die Möglichkeit, über die
Kombination obiger LEDs Diagnose und Fehlercodes auszugeben. Hierbei gilt folgende Codierung:

\begin{itemize}
    \item{Led GREEN blinkt in 1-Sekunden-Perioden:} Der ADuC-7026-Mikrocontroller befindet sich im Bootmodus. Ein Reset des
                                                    ZigBit-Moduls wurde durchgeführt und es wird nun auf Bereitschaft des Moduls
                                                    gewartet.
    \item{Zweimaliges Blinken von YELLOW:} Eine Verbindung mit dem ZigBit-Modul wurde erfolgreich aufgebaut, es wird nun eine 
                                           Initialisierung vorgenommen.
    \item{RED, GREEN leuchten gleichzeitig:} Das ZigBit-Modul wurde erfolgreich initalisiert, auf die Verbindung mit dem 
                                                Funknetzwerk wird gewartet.
    \item{YELLOW leuchtet dauerhaft länger als 10 Sekunden und
          ZIGBIT-STATUS leuchtet dauerhaft:} Es kann keine Verbindung mit dem Funktnetzwerk aufgebaut werden.
    \item{YELLOW leuchtet dauerhaft länger als 10 Sekunden und ZIGBIT-STATUS leuchte nicht:} Es liegt eine Funktionsstörung des
          ZigBit-Moduls vor.
    \item{ZIGBIT-STATUS leuchtet dauerhaft:} Verbindung mit Funknetzwerk wurde verloren oder Powersafe-Modus ist deaktiviert.
    \item{ZIGBIT-STATUS blinkt im 2 Sekunden Interval:} Die Platine arbeitet Ordnungsgemäß, es besteht eine Verbindung mit dem
                                                        Funknetzwerk und der Powersafe-Modus ist aktiv.
\end{itemize}
                                                        
\subsubsection{\emph{MANV-USB-Connector}}
\Abbildungps{!htb}{0.25}{diagramme/Ansteuerung-USB-Stick.pdf}{label}{Datenfluss zwischen Server und Netzwerk}
\Abbildungps{!htb}{0.4}{diagramme/platine_usbstick-crop.pdf}{label}{Platine des USB-Sticks}

Der \emph{MANV-USB-Connector} ist die Schnittstelle zwischen Sensornetz und Computer. Es handelt sich um einen USB-Stick, der einen
ZigBit-Modul beinhaltet. Zusätzlich sind zwei weitere Bauteile enthalten, die das ZigBit-Modul mit Strom versorgen, sowie eine
Umsetzung der UART-Schnittstelle des ZigBit-Moduls auf USB vornehmen. Für die Stromversorgung ist es notwendig, die 5V der
USB-Schnittstelle auf die 3V des ZigBit-Moduls umzusetzen.

\section{Software}
\subsection{MANVNode Firmware}

\subsubsection{Funktionalität}
Die Firmware der \emph{MANVNode} dient dem Entwickeln und Testen eines Treibers für das \emph{ZigBit}-Modul. Dieser
Treiber wird später in den Erste-Hilfe-Sensor integriert werden. Zusätzlich bietet die \emph{MANVNode} Funktionalität,
die zur Simulation des Erste-Hilfe-Sensors gegenüber der \emph{MANVSuite} dient. Hierzu sollen neben dem 
\emph{ZigBit}-Treiber folgende Funktionalitäten integriert werden:

\begin{itemize}
    \item Zufallszahlengenerator zum Senden zufälliger Messwerte
    \item Änderung des Zustands des Zufallszahlengenerators über Tasterkj
    \item Ansteuerung von LEDs zur Ausgabe des Zustands des Zufallszahlengenerators
    \item Ansteuerung eines Piezzo-Summers zur Signalisierung von Alarmen
\end{itemize}


\subsubsection{Ansteuerung des ZigBit-Moduls}
Das \emph{ZigBit}-Modul ist via \emph{UART} mit der \emph{MANVNode} verbunden und wird über das \emph{SerialNet}-Protokoll
angesprochen.

\subsubsection{Ansteuerung des UARTs}
Zur Vermeidung von \emph{Busy-Waiting} soll der Treiber für den \emph{UART}s als Interrupt-Service-Routine (ISR)
realisiert werden. Die Kommunikation zwischen Firmware und \emph{UART}-Treiber soll über Ringpuffer erfolgen:
Zu sendende Daten werden von der Firmware in den Sendepuffer gelegt. Sobald der \emph{UART} sendebereit ist, 
soll der \emph{UART}-Treiber die zu sendenden Daten Zeilenweise aus dem Ringpuffer entnehmen und an
den \emph{UART} übertragen. Analog dazu ist das Empfangen von Daten: Diese werden vom \emph{UART}-Treiber gelesen,
und zeilenweise in den Empfangspuffer gelegt. Die Firmware kann diese nun zeilenweise aus dem Puffer lesen.
Dies entspricht der Semantik einer gepufferten, nicht blockierenden Ein-/Ausgabefunktion.
% Hier Tannenbaum zitieren

\subsubsection{Ansteuerung des ZigBit-Modul}
Die Ansteuerung des \emph{ZigBit}-Moduls soll in folgenden Schritten erfolgen: 

\begin{lstlisting}{frame=entwurfZigBit1}
Uebertragen aktuellen Status
Werte Antwort aus
Arbeite empfangene Befehle ab
Aktiviere den Energiesparmodus
\end{lstlisting}

Diese Schritte können entweder periodisch mittles eines Timer-Interrupts oder als Hauptprogramm in einer Endlosschleife
ausgeführt werden. Wichtig hierbei ist die Berücksichtigung der in Abschnitt~\ref{kommunikation_mit_zigbit_modul}
beschriebenen Synchronisationsprobleme. Hierbei soll wie folgt vorgegangen werden:

\begin{itemize}
    \item{Behandlung asynchron auftretender Ereignisse:} Hierbei ist vor allem das \emph{DATA}-Ereignis wichtig, über das
    der Empfang von Befehlen von der \emph{MANVSuite} signalisiert wird. Alle anderen Events können zunächst vernachlässigt
    werden. Ereignisse, die während des Wartens auf einen Statuscode auftreten, sollen zur späteren Behandlung in einen
    Ringpuffer zwischengespeichert werden. Die Abarbeitung soll dann nach Empfang des Statuscodes erfolgen.

    \item{Behandlung komplexer Rückgabewerte:} Im aktuellen Entwurf werden von der Firmware der \emph{MANVNode} keine 
    \emph{SerialNet}-Befehle verwendet, die eine komplexe Rückgabe haben, so dass dieses Problem eigentlich auftreten 
    sollte. Sollten trotzdem komplexe Rückgaben auftreten, so werden diese wie asynchron auftretende Ereignisse behandelt
    und in den Befehlspuffer gelegt. Bei der späteren Auswertung des Befehlspuffers werden diese dann einfach verworfen.
    So bleibt eine Synchronisation mit dem \emph{ZigBit}-Modul in jedem Fall gewährleistet.
\end{itemize}

Die \emph{SerialNet}-Firmware bietet die Möglichkeit, den Energiesparmodus entweder manuell oder aber in einer 
konfigurierbaren Periode automatisch aufzurufen (Vgl. Abschnitt~\ref{analyse_powermanagement}). Für die
\emph{MANVNode}-Firmware soll der manuelle Aufruf verwendet werden, da dieser deutlich flexibler und einfacher zu
implementieren ist als der automatische Modus: Sobald die Übertragung der Zustandsdaten abgeschlossen ist, und
alle empfangenen Befehle abgearbeitet wurden, soll die Firmware über den Befehl \emph{AT+WSLEEP} den Energiesparmodus
aufrufen. Hierdurch ist gewährleistet, dass das \emph{ZigBit}-Modul nur so lange wie zur Bearbeitung aller Aufgaben
notwendig in Betrieb ist. Nachdem das \emph{ZigBit}-Modul aus dem Energiesparmodus zurückkehrt wird die Abarbeitung
der Ansteuerungsroutine von vorne begonnen. Das Aufwachen aus dem Energiesparmodus wird der \emph{MANVNode}-Firmware
über eine Zustandsänderung des \emph{CTS}-Registers des \emph{UART} signalisiert. Zusätzlich wird der \emph{AT+WSLEEP}
Befehl mit dem Statuscode \emph{OK} quittiert, sobald das \emph{ZigBit}-Modul wieder betriebsbereit ist. 


Insgesamt ergibt sich für die Ansteuerung des ZigBit Moduls folgender Ablauf:

\begin{lstlisting}{frame=entwurfZigBit2}
Uebertrage aktuellen Status
Wiederhole:
    Antwort von ZigBit-Modul lesen
    Falls Antwort != Statusnachricht:
        Antwort in Befehlspuffer legen
bis Antwort == Statusnachricht    
Arbeite Befehlspuffer ab
Aktiviere den Energiesparmodus
\end{lstlisting}

Dieser Ablauf ist in Abbildung~\ref{zigbit_mainloop} als Ablaufdiagramm dargestellt.
\Abbildungps{!htb}{0.5}{diagramme/Flussdiagramm-Firmware-Mainloop.pdf}{zigbit_mainloop}{Ablauf der Ansteuerung des \emph{ZigBit}-Moduls.}

\subsection{MANVSuite}

In dieser Arbeit wurde ein Java-Treiber (MANVConnector) entworfen und implemtiert. Dieser Treiber realisiert die die 
Anbindung an die von Herrn Tepelmann in \cite{Jan} entwickelte MANVSuite. 


\subsection{MANVConnector}
\Abbildungps{!htb}{1.25}{diagramme/manv_connector_sequenzdiagramm_main.pdf}{sequenzdiagramm_manvconnector}{Interaktion der einzelnen Threads des \emphf{MANVConnectors}}


Der \emph{MANVConnector} hat einerseits die Aufgabe, Daten die von dem \emph{MANV-USB-Connector} empfangen wurden in Corba-Events 
umzusetzen, und an den \emph{MANVServer} weiterzuleiten. Andererseits empfängt sie Corba-Events vom \emph{MANVServer}, dekodiert
diese und sendet diese in Form von Sensornetz Befehlen an die zuständigen \emph{MANVNodes} weiter.

Bei der Kommunikation mit dem auf dem USB-Stick aufgebrachten ZigBee-Modul stellen sich grundsätzlich die selben 
Synchronisierungsprobleme wie in der \emph{MANVFirmware}. Da der \emph{MANVConnector} jedoch alle Möglichkeiten der Java-Virtual-Machine
nutzen kann, lassen sich diese deutlich einfacher und eleganter lösen. 

In der MANVFirmware werden hierzu mehrere Ringpuffer verwendet, welche die empfangenen Daten speichern. Diese finden sich auch
im MANVConnector wieder. Hierbei werden an Stelle von Ringpuffern priorisierte Warteschlangen, sogenannte Queues verwendet. 

\begin{itemize}
    \item{Die Command Queue:} In dieser Queue werden alle zu sendenden Befehle gespeichert.
    \item{Die Event Queue:} In dieser Queue werden alle empfangenen Ereignisse gespeichert.
    \item{Die Result Queue:} In dieser Queue werden alle bereits gesendeten Befehle zusammen mit dem
                             Resultat, das dieser Befehl hatte, gespeichert.
\end{itemize}

Jedes Element der einzelnen Queues verfügt über eine Priorität; die Queues sorgen dafür, dass der
Zugriff nach Priotität sortiert erfolgt. Hierdurch wird sicher gestellt, dass wichtige Ereignisse
wie z.B. ein Alarm bevorzugt ausgeliefert werden.

Der MANVConnector verfügt neben dem Hauptthread über drei weitere Threads:

\begin{itemize}
    \item{SocketWriter:} Dieser Thread entnimmt Befehle aus der \textsl{CommandQueue} und sendet
                         diese über den MANV-USB-Connector an das Sensornetz. Nun blockiert der
                         Thread so lange, bis das Ergebnis des Befehls zur Verfügung steht.
                         Sobald dies der Fall ist, wird der Befehl zusammen mit dem Ergebnis in
                         die \textsl{ResultQueue} eingefügt.
    \item{SocketReader:} Dieser Thread empfängt Daten aus dem Sensornetz. Handelt es sich um ein
                         Ergebnis, so wird dies dem \textsl{SocketWrite} signalisiert, und das
                         Ergebnis zur Abholung durch den \textsl{SocketWriter} zur Verfügung
                         gestellt. Handelt es sich hingegen um einen Ereignis, so wird dieses
                         in die \textsl{EventQueue} eingefügt.
    \item{CorbaSender:} Dieser Thread ist für die Kommunikation mit MANVServer zuständig.
                         Befehle, die vom MANVServer empfangen werden, werden für das 
                         Sensornetz aufbereitet und in die CommandQueue eingestellt.
                         Ausserdem werden Ereignisse aus der \textsl{EventQueue} entnommen,
                         in Corba-Events übersetzt und an den MANVServer zugestellt.
\end{itemize}                          

Die Kommunikation der Threads untereinander findet komplett über die oben genannten Warteschlangen statt. Das Zusammenwirken der Threads
ist in Abbildung~\ref{sequenzdiagramm_manvconnector} dargestellt:\\

Zunächst erzeugt der Hauptthread (\emph{MANVConnector.main}) die drei Warteschlangen. Dannach werden die beiden Threads
\emph{SocketWriter} und \emph{SocketReader} erstellt. Dannach ruft der Hauptthread die Corba-Inizialisierungsroutine auf und startet
den \emph{CorbaSender}-Thread.\\

Beim Starten der Threads, wird sichergestellt, dass alle Zugriff auf alle drei Warteschlangen haben. Die Konsequente Verwendung
der Warteschlangen zur Kommunikation der Threads untereinander löst bereits viele Synchronisationsprobleme. Grundsätzlich findet
Kommunikation zwischen den Threads rein über die Warteschlangen statt. Einzige Ausnahme ist das Abrufen des Ergebnisses eines
gesendeten Kommandos. 

\subsubsection{Synchronisation der Befehlsübertragung}


%%%-----------------------------------------------UMSETZUNG/VORBEREITUNGEN-------------------------------------------------------------
\chapter{Praktische Realisierung des Sensornetzes}\label{Implementierung}

In diesem Kapitel wird die Implementierung des Sensornetzes, dabei auftretende Probleme sowie deren Lösungen, beschreiben.

\section{Hardware}
Nachfolgend wird die praktische Umsetzung der zwei verschiedenen Hardwarekomponenten, also des \emph{MANV-USB-Sticks} und der
\emph{MANVNode} beschrieben. In beiden Fällen musste hierzu ein \emph{ZigBit}-Modul mit einer anderen 
Hardwarekomponente verbunden werden. Dies kann grundsätzlich entweder über \emph{SPI} oder \emph{UART} erfolgen.
Bei \emph{SPI} handelt es sich um einen seriellen Bus, der besonders im Bereich der Mikrocontroller verbreitet ist,
und zur Anbindung von weiterer Peripherie dient. \emph{UART} war hingegen vor allem als serielle Schnittstellen
an \emph{PCs} verbreitet, und wurde hauptsächlich für die Anbindung von Modems verwendet. Heutzutage wurde diese
Schnittstelle weitestgehend von \emph{USB} verdrängt. Beim Einsatz der \emph{SerialNet}-Firmware scheidet die
Anbindung über \emph{SPI} jedoch aus, da dies von der Firmware nicht unterstützt wird.

\subsection{Anbindung des ZigBit-Moduls an den USB-Port}
\nomenclature{IC}{Integrierter Schaltkreis}
Zur Realisierung des \emph{USB-Sticks} bietet sich die Verwendung des \emph{UARTs} an, da eine breite Palette 
von fertigen \emph{ICs} zur Verfügung steht, mit deren Hilfe ein \emph{UART} an einen \emph{USB-Bus} angebunden werden 
kann. Auf Seiten des Betriebssystems präsentiert sich der \emph{IC} wie ein serieller Port und kann softwareseitig
wie ein solcher angesprochen werden. Die Firma \emph{Embedded Projects} bietet mit der \emph{USB-UART-Bridge}
eine fertige Lösung an, auf der ein entsprechender \emph{IC}, ein Spannungsregler für 3.3~V-Versorgungsspannung
sowie ein \emph{USB-Anschluss} bereits integriert sind. Das \emph{ZigBit}-Modul wird nun auf die in 
Abbildung~\ref{platine_usbstick} dargestellte Platine aufgelötet. Diese Platine kann nun selbst mit Hilfe einer
Stiftleiste huckepack auf die \emph{USB-UART-Bridge} aufgelötet werden, um die nötige Formstabilität zu erhalten.

\subsection{Herstellung und Bestückung der Platine der \emph{MANVNode}}
\Abbildungps{!htb}{0.75}{diagramme/usbstick_makro.png}{usbstick}{\emph{MANV-USB-Stick}.}
\Abbildungps{!htb}{1}{diagramme/manvnode.jpg}{manvnode}{Fertig bestückte Platine der \emph{MANVNode.}}
\nomenclature{nF}{Nanofarad}
\nomenclature{Flashen}{Programmieren eines Flashspeichers.}
\nomenclature{Lötstopplack}{Auf Leiterplatten aufgetragene Lackschicht, die die Leiterbahnen schützt.}
Auch bei der \emph{MANVNode} erfolgt eine Anbindung des \emph{ZigBit}-Moduls über \emph{UART}. Hierbei muss
beachtet werden, dass der \emph{UART} auch für das Flashen des Microcontrollers benötigt wird. Daher verfügt
die Platine über eine Reihe von Jumpern, mit denen das \emph{ZigBit}-Modul vom \emph{UART} abgetrennt werden
kann.\\
        \\
Der \emph{ADuC-7026-Mikrocontroller} benötigt drei 470~nF-Kondensatoren zur Stabilisierung der
Stromversorgung. Diese waren in den ersten beiden Versionen der Platine zu weit vom \emph{Mikrocontroller} entfernt 
aufgebracht worden, was dazu führte, dass der \emph{Mikrocontroller} beim Flashvorgang, dazu neigte, spontan 
abzustürzen. Dieser Fehler konnte durch die Verringerung des Abstandes zwischen Mikrocontroller und 
Kondensatoren behoben werden.


\section{Firmware}
Die Firmware des Erste-Hilfe-Sensors wurde um einen Treiber für das \emph{ZigBit}-Modul ergänzt. Die \emph{Firmware}
wurde in der Programmiersprache \emph{C} geschrieben, und setzt direkt auf die Hardware des \emph{ADuC} auf. Es wurde 
lediglich die von \emph{Rowley Crossworks} angebotene Standardbibliothek verwendet, die einige praktische Funktionen wie
Stringmanipulation, einen Interrupthandler und einen fertigen \emph{Startup-Code} bietet.
Die Entwicklung von Software für einen Microcontroller zeichnet sich durch die Abwesenheit eines Betriebssystems 
aus. Ein Großteil der Funktionalität, die man von der Entwicklung von Software für einen gewöhnlichen Mikrocomputer 
gewohnt ist, ist schlichtweg nicht vorhanden. Hier sind insbesondere eine automatische Speicherverwaltung sowie
\emph{Threads} und \emph{Prozesse} zu \emph{erwähnen}. Da der \emph{Erste-Hilfe-Sensor} viele Aufgaben gleichzeitig 
erfüllen muss, stellt dies eine ernst zu nehmende Herausforderung dar. Das Problem wurde durch ein
\emph{Interrupt}-getriebenes Programmiermodell gelöst. Zu beachten sind hierbei die sehr beschränkten Ressourcen
des \emph{Mikrocontrollers}. Insbesondere der Speicher ist mit 16~kB sehr knapp bemessen. 

\subsection{Verwendete Interrupts}
\nomenclature{Crossworks}{Entwicklungsumgebung für \emph{Mikrocontroller-Firmware}.}
\nomenclature{Interrupt}{Unterbrechungsaufforderung. Die Abarbeitung des aktuellen Programmcodes wird sofort beendet,
                         und eine \emph{ISR} wird ausgeführt.}
\nomenclature{Startup-Code}{Programmcode, der die Initialisierung eines \emph{Mikrocontrollers} durchführt.}
Es werden folgende \emph{Interrupts} verwendet:

\emph{Timer0}: Dieser \emph{Timer-Interrupt} führt die Patientenüberwachung durch. Die einzelnen Sensoren werden abgefragt,
und eine Analyse der empfangenen Daten wird durchgeführt.

\nomenclature{Piezo-Summer}{Elektrisches Bauteil, dass einen Summton generiert, sobald man eine Spannung anlegt.}
\emph{Timer1}: Dieser \emph{Interrupt} führt einige periodische Aufgaben durch. Zunächst werden die am Sensor vorhandenen
Taster abgefragt (Alarm stumm schalten, Alarm manuell auslösen etc.). Danach wird überprüft, in welchem Zustand
der Sensor sich aktuell befindet, also z.B. ob ein Alarm aufgetreten ist oder ob der Patient sich in einem
guten Zustand befindet. Abhängig hiervon werden nun \emph{LEDs} und ein angeschlossener \emph{Piezo}-Summer geschaltet,
um den Zustand nach außen zu signalisieren. Zuletzt wird noch überprüft, wann zuletzt eine Übertragung
des Zustands des Sensors an die \emph{MANVSuite} erfolgt ist. Liegt dies länger als ein konfiguriertes Zeitintervall
zurück, so wird eine Übertragung des aktuellen Zustands veranlasst. Ein sinnvolle Länge für dieses Zeitintervall
ist ein Wert von 5-10~s.

\subsubsection{Zugriff auf UART}
Der \emph{ADuC} verfügt über einen \emph{UART-Interrupt}, welcher eine Statusänderung des \emph{UARTs} signalisiert. 
Im Register \emph{COMIEN0} wird konfiguriert, welche Zustände über den \emph{Interrupt} signalisiert werden sollen. Tritt
nun einer dieser Zustände auf, so wird der \emph{UART-Interrupt} ausgelöst. In der \emph{ISR} muss nun überprüft
werden, welches Ereignis zum Auslösen des \emph{Interrupts} geführt hat. Diese Information ist im Register \emph{COMSTA0} 
gespeichert. Wichtig ist an dieser Stelle, dass auch mehrere Ereignisse gleichzeitig auftreten können. 
Dies muss in der \emph{ISR} berücksichtigt werden, da sonst Ereignisse verloren gehen können.

Das eigentliche Senden und Empfangen erfolgt über die beiden Register \emph{COMRX} und \emph{COMTX}.
Zum Senden wird hierzu ein einzelnes Zeichen in \emph{COMTX} gelegt. Nun muss eine gewisse Zeit gewartet
werden, bis das Zeichen gesendet wurde und das nächste Zeichen in \emph{COMTX} gelegt werden kann.
Für das Empfangen wird das Register \emph{COMRX} in analoger Weise verwendet. Ob das nächste
Zeichen empfangen bzw. gesendet werden kann, kann mit Hilfe der Bits \emph{DR} (\emph{"`Data Ready"'} - Daten liegen vor)
bzw. \emph{TEMT} (\emph{"`Transmit Buffer empty"'} - Daten können gesendet werden) bestimmt werden.

Die einfachste Methode wäre hierbei, in einer Schleife \emph{Busy-Waiting} zu betreiben und so lange zu warten,
bis sich eins der beiden Bits verändert. Dies wäre jedoch sehr aufwendig und würde den \emph{Mikrocontroller}
unnötig lange blockieren. Statt dessen wird der Zustand der beiden Register nur dann überprüft, wenn ein 
\emph{UART-Interrupt} aufgetreten ist.

Zum Senden und Empfangen von Daten werden zwei Ringpuffer verwendet. Möchte ein Unterprogramm Daten senden,
so greift es nicht direkt auf die \emph{UART}-Schnittstelle zu, sondern legt diese Daten lediglich in den Sendepuffer.
Das eigentliche Senden wird nun vom \emph{UART-Interrupthandler} durchgeführt; das Unterprogramm kann weiter arbeiten,
ohne auf das fertige Senden der Daten warten zu müssen.

\nomenclature{Idle Task}{Aufgabe, die immer dann bearbeitet wird, wenn keine anderen Aufgaben anstehen.}
Das Empfangen von Daten erfolgt analog. Jedesmal, wenn ein Zeichen von der seriellen Schnitstelle empfangen
wurde, wird dieses in den Empfangspuffer gelegt. Das Abarbeiten des Empfangspuffers erfolgt nun als \emph{Idle-Task}:
Immer dann, wenn der \emph{Mikrocontroller} gerade keine anderen Aufgaben erfüllen muss, wird der Empfangspuffer
abgearbeitet und eventuell empfangene Befehle werden ausgeführt. Dies kann natürlich jederzeit durch 
die Behandlung von \emph{Interrupts} unterbrochen werden.

\section{MANVConnector}

Die folgenden Abschnitte beschreiben, wie der Zugriff vom \emph{MANVConnector} auf den \emph{MANV-USB-Stick} 
konkret erfolgt. Der \emph{MANV-USB-Stick} verhält sich hierbei dem Betriebssystem gegenüber wie eine serielle 
Schnittstelle.

\nomenclature{Socket}{Softwareabstraktion einer Netzwerkverbindung die sich verhält wie ein Dateizugriff.}
\subsection{Zugriff auf die serielle Schnittstelle}
Der Zugriff auf die serielle Schnittstelle mittels Java bringt einige praktische Probleme mit sich.
Der \emph{Java}-Standard bietet von Haus aus keine Möglichkeit, auf die serielle Schnittstelle zuzugreifen.
Zwar gibt es einzelne Lösungen von Dritt\-an\-bie\-tern wie z.B. \emph{RxTx} oder \emph{JavaComm}, allerdings
sind diese meist plattformabhängig und funktionieren meist nur auf wenigen Betriebssystemen. Da die
in dieser Diplomarbeit entworfene Lösung auf möglichst vielen Betriebssystemen einsetzbar sein soll, 
wurde entschieden, den Zugriff auf die serielle Schnittstelle stattdessen in einem austauschbaren
Modul zu kapseln. Der Zugriff auf das Modul erfolgt über eine \emph{Socketverbindung}, welche in \emph{Java}
standardmäßig verfügbar ist. Das Modul kann nun in jeder beliebigen Programmiersprache implementiert
werden. Es ist z.B. möglich, für verschiedene Betriebssysteme verschiedene Module anzubieten.
Für die Zwecke dieser Diplomarbeit wurde beispielhaft ein Modul für \emph{UNIX} in der Programmiersprache
Python implementiert, welche unter vielen \emph{UNIX}-artigen Betriebssystemen wie \emph{Linux}, \emph{MacOS}, 
\emph{BSD} und co.  bereits zur Standardinstallation gehört. Lediglich das \emph{PySerial}-Modul muss nachinstalliert werden.
% TODO: Hier Zitat des PySerial Moduls

\begin{lstlisting}{Seriellen Port auf Socket umleiten}
#!/usr/bin/env python -OO
# MANV-Sensor - Serial to socket
# for Java to ZigBit connector 

PORT = 4711 ## Port fuer Socketverbindung.
import serial, socket, select, sys, SocketServer, time

class MANVSerialToSocketHandler(SocketServer.\
        BaseRequestHandler):

    def handle(self):
        ## Handler-Hook, in den der SocketServer einsteigt
        print "Connection to socket opened"
        self.openSerial()

        print "Serial connected"
        self.tunnel()

    def openSerial(self):
        ## Serielle Schnittstelle oeffnen.
        ## Welche das genau ist, wird auf der Kommandozeile  
        ## als Parameter uebergeben.
        self.serial = serial.Serial(sys.argv[1], 38400, 
                                    timeout=1)

    def shutdown(self):
        ## Socket-Verbidung wird geschlossen -> aufraeumen
        self.serial.close()

    def tunnel(self):
        ## Weiterleiten von Daten:
        ## Socket -> Serial
        ## Serial -> Socket
        end = False
        print "Entering tunnel mode"

        while not end:
            ## Select verwenden, um socket und seriellen  
            ## Port ressourcensparend zu ueberwachen. 
            ## Funktioniert leider nur unter UNIX.
            readylist = select.select((self.request, 
                                       self.serial),
                                       (), ())

            ## Behandeln aller empfangenen Daten
            for readysocket in readylist[0]:

                ## Fallunterscheidung zwischen seriellem 
                ## Port und Socket, da hier die 
                ## Semantik unterschiedlich ist.
                if readysocket == self.serial:

                    ## Daten  wurden auf seriellem Port 
                    ## empfangen -> auf socket schreiben.
                    data = self.serial.read(1024)
                    self.request.send(data)

                elif readysocket == self.request:

                    ## Daten wurden auf Socket empfangen
                    ## -> auf seriellen Port schreiben
                    data = self.request.recv(1024)

                    if data:
                        self.serial.write(data)

                    else:
                        ## EOF empfangen -> shutdown
                        end = True

                ## select hat etwas zurueck geliefert, das 
                ## weder Socket noch serieller Port ist
                ## -> shutdown
                else:
                    data = ""
                    end = True

                ## Daten zu Debuggingzwecken auf Konsole 
                ## ausgeben.
                print data

        ## Aufraeumen: Seriellen Port schliessen 
        self.serial.close()
        self.request.close()
        print "Leaving tunnel mode"


class ReusableServer(SocketServer.TCPServer):

    def server_bind(self):
        ## Socket wiederverwendbar machen.
        ## Wenn man das nicht tut, ist er nach Beenden 
        ## des Programms fuer 5 Minuten nicht ansprechbar.
        self.socket.setsockopt(socket.SOL_SOCKET,
                               socket.SO_REUSEADDR, 1)
        self.socket.bind(self.server_address)

## Einstiegspunkt 
if __name__ == '__main__':
      server = ReusableServer(('localhost', PORT), 
                              MANVSerialToSocketHandler)
      server.serve_forever()
\end{lstlisting}

\nomenclature{Port}{Schnittstelle}
Für die Überwachung von \emph{Socket} und seriellem \emph{Port} wird in diesem Programm der Systemaufruf
\emph{select()} verwendet. Wird dieser aufgerufen, gibt das Programm so lange den \emph{Prozessor}
für andere Programme frei, bis Daten empfangen wurden. Sobald dies der Fall ist, kehrt
der \emph{select()}-Aufruf zurück und liefert eine Liste mit allen Quellen, von denen Daten
empfangen wurden. Diese können nun mit den entsprechenden Aufrufen von \emph{recv()} 
(im Falle von \emph{Sockets}) bzw. \emph{read()} (für den seriellen \emph{Port}) abgeholt werden.
Diese Lösung funktioniert jedoch in dieser Form nur unter \emph{UNIX}. Unter \emph{Windows} steht der 
\emph{select()}-Aufruf zwar zur Verfügung, jedoch ist es hier nicht möglich, diesen zur
Überwachung des seriellen \emph{Ports} zu verwenden. Das obige Programm lässt sich hier jedoch
recht einfach durch die Verwendung von \emph{Threads} anpassen, indem für jede Richtung (also
Seriell$\,\rightarrow{}\,$\emph{Socket} und \emph{Socket}$\,\leftarrow{}\,$Seriell) jeweils ein eigener \emph{Thread} verwendet
wird. Alternativ empfiehlt sich eine Implementierung mit Hilfe der Programmiersprache \emph{C++}.

\subsection{Automatische Erkennung des \emph{ZigBit-USB-Sticks}}\label{usb_stick_erkennen}
\nomenclature{Bus}{Verbindung zwischen mehreren Geräten, bei denen alle Geräte an einen gemeinsamen Übertragungsweg
                   angeschlossen sind.}  
Bei der Verwendung von \emph{USB}-Geräten stellt sich bei vielen Betriebssystemen das Problem,
dass die Benennung bei jedem Einstecken des Gerätes anders sein kann. Dies kann insbesondere
dann auftreten, wenn mehrere gleichartige Geräte am \emph{USB-Bus} (z.B. mehrere 
\emph{USB-Serial}-Geräte) vorhanden sind. Um nicht nach jedem Einsteckvorgang oder Neustart des Betriebssystems die
Konfiguration des \emph{MANV-Connectors} anpassen zu müssen, wird der Name des \emph{USB}-Sticks
nicht fest konfiguriert, sondern stattdessen eine automatische Erkennung durchgeführt.
Unter \emph{Linux} stehen mit \emph{libusb} und dem \emph{/sys}-Dateisystem zwei geeignete Werkzeuge
zur Verfügung, um diese Erkennung durchzuführen.\\
Zunächst werden mit Hilfe der \emph{libusb} alle \emph{USB-Busse} nach dem auf dem \emph{ZigBit}-\emph{USB-Stick}
verwendeten \emph{IC} vom Typ \emph{cp210x} durchsucht. Dieser hat die Herstellerkennung \emph{0x10C4} und die
Produktkennung \emph{0xEA60}. Unter Python steht mit \emph{python-usb} eine geeignete Schnittstelle zu
\emph{libusb} zur Verfügung.

\begin{lstlisting}{USB-Bus durchsuchen}
import usb
def searchUSB(self):

    ## Durchsuchen alle USB-Busse 
    for bus in usb.busses():

        ## Durchsuchen aller Geraete dieses Busses
        for dev in bus.devices:

            ## Betrachten der Vendor-ID und der Produkt-ID
            if dev.idVendor == 0x10c4 and \
               dev.idProduct == 0xea60:
               
                ## Ein entsprechendes Geraet wurde 
                ## gefunden wir geben dieses 
                ## zurueck 
                return (bus.dirname, dev.filename)
   
    print "No devices found on USB-Bus."
    return False
\end{lstlisting}

Wenn am \emph{USB-Bus} kein entsprechendes Gerät gefunden wurde, kann bereits jetzt die Suche abgebrochen werden.
Wird jedoch ein passendes Gerät gefunden, müssen zunächst alle passenden \emph{Device-Nodes}\footnote{Unter UNIX
verhalten sich Hardwaregeräte grundsätzlich wie Dateien und finden sich daher im Dateisystem unterhalb
des Verzeichnisses \emph{/dev} als virtuelle Datei wieder. Eine solche Datei wird als \emph{Device-Node} bezeichnet.} im
\emph{/dev}-Verzeichnis gefunden werden. Hierzu bietet sich eine Suche im \emph{/sys}-Dateisystem an. Hierfür sind keine
externen Module notwendig, es können die normalen Dateisystem-Operationen des \emph{OS}-Moduls verwendet werden:

\begin{lstlisting}{Devicenode lokalisieren}
import os
def searchSYS(self):

    ## Suchen aller usb-serial-Geraete ueber /sys 
    devices = []

    for device in \
        os.listdir('/sys/bus/usb-serial/devices'):
        ## Alle Geraete in Rueckgabe einfuegen
        devices.append(device)

    return devices
\end{lstlisting}

\nomenclature{Debuggen}{Suchen und Entfernen von Fehlern.}
Da an den \emph{USB-Bus} beliebige \emph{USB-Serial}-Geräte angeschlossen sein können\footnote{Während der Entwicklung
was es nicht ungewöhnlich, dass bis zu 3 Geräte angeschlossen waren: Bei einem handelte es sich um den
\emph{ZigBit-USB-Stick}, zwei weitere dienten zum \emph{Flashen} und \emph{Debuggen} der angeschlossenen 
\emph{Mikrocontroller}.}, 
muss nun für jedes der angeschlossenen Geräte überprüft werden, ob es sich um einen \emph{ZigBit}-\emph{USB}-Stick
handelt. Hierzu wird an das Gerät ein Befehl gesendet, der zur Abfrage der Versionsnummer der verwendeten
\emph{Serialnet}-Firmware dient. Die Rückgabe wird nun mit dem erwarteten Wert verglichen. Hierzu wird das
\emph{python-serial}-Modul verwendet:

\begin{lstlisting}{ZigBit-Modul erkennen}
import serial
def testDevice(self, deviceName):
    ret = []

    ## Versuch, eine Verbindung aufzubauen und die 
    ## Versionskennung der Serialnet-Firmware abzufragen. 
    try:
        ser = serial.Serial("/dev/%s" % deviceName, 38400,\
                            timeout=3)
        ser.write("ATI\r\n") ## ATI - Befehl zum Anfragen
                             ## der Versionsnummer
        out = ser.read(74)
        
        ## Wurde ueberhaupt eine Antwort empfangen? 
        if out:

            ## Antwort in Bestandteile zerlegen
            formatedOut = out.split("\r\n")

            ## Hat die Antwort das richtige Format?
            if len(formatedOut) == 5:
                vendor  = formatedOut[1]
                name    = formatedOut[2]
                version = formatedOut[3]
                mac     = formatedOut[4]

                ## Herstellerkennung und Modell ueberprufen
                if vendor == 'ATMEL' and name == 'ZIGBIT':

                    ## Ein ZigBit-Modul wurde erkannt
                    return (vendor, name, version, mac)
                else:

                    ## Es wurde kein ZigBit-Modul erkannt
                    return False
            else:
                ## Antwort hatte falsches Format -> Es 
                ## handelt sich nicht um ein 
                ## ZigBit-Modul
                return False

        else:
            ## Keine Antwort empfangen
            return False

    ## Fehler beim Zugriff auf serielle Schnittstelle
    except serial.SerialException:
        return False
\end{lstlisting}

\subsection{CORBA-Anbindung}
\Abbildungps{!htb}{0.5}{diagramme/corba.pdf}{corba}{Komponentendiagramm der \emph{CORBA}-Schnittstellen. (Quelle: Jan Tepelmann)}
\nomenclature{Nameserver}{Softwaredienst, der eine Übersetzung von Namen in Adressen vornimmt.}
Die \emph{CORBA}-Anbindung des \emph{MANVConnectors} an den \emph{MANVServer} erfolgt über
die in \cite{Jan} spezifizierten \emph{Interfaces}. Diese müssen über den \emph{Java Classpath}
eingebunden werden. Dies geschieht über den Parameter \emph{-cp} auf der Kommandozeile
der \emph{JVM} bzw. des \emph{Java-Compilers}. Der \emph{Connector} schickt über die 
\emph{Server\_Incoming}-Schnittstelle eingehende Daten und Ereignisse an den \emph{MANVServer}.
Damit der \emph{MANV\-Server} umgekehrt auch Befehle an den \emph{MANVConnector} schicken kann,
muss dieser selbst die \emph{Connector\_Commands}-Schnittstelle implementieren. Dies
geschieht in der Klasse \emph{CommandsImpl}\footnote{Eine
genaue Beschreibung der einzelnen Dateien, Klassen und Funktionalitäten findet sich in
Anhang~\label{anhang_beschreibung_software}.}.\\

Zur Lokalisierung der einzelnen \emph{CORBA}-Schnittstellen dient der \emph{Nameserver}.
Die Adresse des Nameservers muss im Voraus bekannt sein. Meist läuft dieser auf dem selben Rechner
wie der \emph{MANVServer}. Die Adresse wird beim Start des \emph{MANVConnectors} als 
Kommandozeilenparameter "`--ORBInitRef NameService"' übergeben\footnote{Wenn z.B. Nameserver, \emph{MANVServer}
und \emph{MANVConnector} auf dem selben Rechner laufen, so lautet der entsprechende Parameter:
"`--ORBInitRef NameService=corbaloc::localhost:1771/NameService"'.}. Die Verbindung zum \emph{Nameserver}
wird nun beim Start des \emph{MANVConnectors} in der Methode \emph{initCORBA()} aufgebaut. 
Sobald diese hergestellt wird, wird nach der Adresse des \emph{MANVServers} gefragt. Dies geschieht über
eine Anfrage nach den Schnittstellen des Servers, also \emph{Server\_Incoming} und \emph{Server\_Control}. 
Zuletzt muss sich der \emph{MANVConnector} noch beim \emph{MANVServer} registrieren, damit dieser
weiß, an wen er zu sendende Befehle zustellen muss. Hierzu wird die Methode \emph{registerConnector()}
der \emph{Server\_Control}-Schnittstelle verwendet. Nach diesem Schritt ist die Initialisierung der
\emph{CORBA}-Anbindung abgeschlossen, und der \emph{CorbaSender}-Thread, der eingehende Ereignisse
an den \emph{MANVServer} sendet, kann gestartet werden.



%%%------------------------------------------------ERGEBNISSE(KEINE BEWERTUNG)----------------------------------------------------

\chapter{Ergebnisse}
\section{Leistungsaufnahme}
\Abbildungps{!htb}{0.75}{oszi/normalbetrieb_router.jpg}{oszi_normalbetrieb_router}{Router oder Koordinator im Normalbetrieb. 
    Der Stromverbrauch beträgt hierbei konstant 23,4mA, was ca. 70mW entspricht.} 
\Abbildungps{!htb}{0.75}{oszi/normalbetrieb_client.jpg}{oszi_normalbetrieb_client}{Endknotens im Normalbetrieb ohne 
    Energiesparmodus. Auffällig sind die Peaks, in denen sich der Knoten im Empfangsmodus befindet. Die Peaks haben eine 
    Länge von 40ms und einen Betrag von 23,4mA. Ausserhalb eines Peaks beträgt der Stromverbrauch 10,7mA. Insgesamt ergibt 
    sich ein Mittelwert von 11,18mA, was ca. 33,54mW entspricht. Dies ist weniger als die Hälfte des Stromverbrauchs
    eines Routers oder Koordinators.} 
\Abbildungps{!htb}{0.75}{oszi/empfangen_details.jpg}{oszi_empfangen_details}{Detailaufnahme eines Empfangsmodus-Peaks.
    Erkennbar ist die Impulsbreite von 40ms.}
\Abbildungps{!htb}{0.75}{oszi/not_joined.jpg}{oszi_not_joined}{Bei diesem Client wurde der Empfangsmodus deaktiviert. Der 
    Stromverbrauch beträgt konstant 10,7mA.}
\Abbildungps{!htb}{0.75}{oszi/joined.jpg}{oszi_joined}{Beitrittsvorgang eines Clients in ein Netzwerk. Bevor der Client dem 
    Netzwerk beigetreten ist, befindet er sich durchgehend im Empfangsmodus und benötigt 23,4mA. Sobald er dem Netzwerk
    beigetreten ist, entspricht entspricht der Stromverbrauch dem in Abbildung~\ref{oszi_normalbetrieb_client} dargestellten
    Verhalten.} 
\Abbildungps{!htb}{0.75}{oszi/joined2.jpg}{oszi_connection_lost}{Dieser Client hat kurzzeitig die Verbindung zum Netzwerk 
    verloren. In dieser Zeit befand er sich im Dauerempfangsmodus und hat 23,4mA. benötigt.}
\Abbildungps{!htb}{0.75}{oszi/erwachen_1.jpg}{oszi_erwachen}{Dieser Client ist gerade aus dem Energiesparmodus erwacht. 
    Erkennbar ist, dass während der Schlafphase der Stromverbrauch praktisch auf 0 sinkt. Direkt nach dem Erwachen
    wird ein Empfangsvorgang durchgeführt; die danach folgende Phase ist analgo zu Abbildung~\ref{oszi_normalbetrieb_client}.
    Gut zu Erkennen ist, dass der Energiesparmodus praktisch keinen zusätzlichen Overhead bedeutet.}
\Abbildungps{!htb}{0.75}{oszi/powersave_periodic.jpg}{oszi_powersave_periodic}{Dieser Client nutzt das periodische Aufrufen
    des Energiesparmodus. Die Schlafphase hat hier eine Länge von 5 Sekunden, die Wachphase ist genau 1 Sekunde lang. Der
    Stromverbrauch beträgt im Mittel nur noch 1,86mA, was ca. 5,58mW entspricht.}

\section{Reichweite}

%%%------------------------------------------------DISKUSSION(BEWERTUNG DER ERGEBNISSE)---------------------------------------------

\chapter{Diskussion}\label{Diskussion}
In Abschnitt~\ref{Ergebnisse} wurde die Leistungsfähigkeit des Funknetzwerkes untersucht.
Es wurde gezeigt, dass sowohl Reichweite als auch Leistungsaufnahme die Anforderungen
dieser Arbeit erfüllen. Auch wurde die Erweiterbarkeit des Netzwerkes durch \emph{Router} gezeigt.\\
\\
Darüber hinaus wurde gezeigt, dass eine Verbindung des verwendeten \emph{ZigBee}-Modul
mit dem \emph{ADuC}-Mikrocontroller möglich ist, und dass das so entstehende Gesamtsystem
in der Lage ist, die zu überwachenden Vitaldaten zuverlässig an die Basistation zu übermitteln.
Es wurde ebenfalls gezeigt, dass es möglich ist, Befehle über das \emph{ZigBee}-Netzwerk
an den Mikrocontroller zu senden. Da der für diese Tests eingesetzte \emph{ADuC-7026}-Mikrocontroller
der selben Familie wie der in dem \emph{Erste-Hilfe-Sensor} verwendete \emph{Mikrocontroller} handelt,
können diese Ergebnisse ohne weiteres auf die Hardware des \emph{Erste-Hilfe-Sensors} übertragen werden; 
die Integrierbarkeit in diesen konnte somit gezeigt werden.\\
Durch die erhöhte Leistungsaufnahme, insbesondere bei Lastspitzen (\emph{Peaks}), muss eine Anpassung
der Stromversorgung des \emph{Erste-Hilfe-Sensors} erfolgen, wenn das \emph{ZigBee}-Modul integriert
werden soll. Die genaue Analyse auftretender Lastspitzen findet sich in Abschnitt~\ref{leistungsaufnahme_normalbetrieb}.
Insbesondere sollten hierbei auch die Ergebnisse in Abschnitt~\ref{leistungsaufnahme_sonderfaelle} beachtet 
werden, die auftretende Sonderfälle beschreiben.\\
\\
Softwareseitig konnte gezeigt werden, dass eine vollständige Kapselung aller hardwarespezifischen Teile
in einer austauschbaren Komponente möglich ist. Dies bietet den Vorteil, dass bei einem Wechsel 
der Technologie des Sensornetzwerks lediglich der \emph{MANVConnector} ausgewechselt werden muss,
alle anderen Teile können ohne Änderung weiterverwendet werden.\\
\\
Durch die in Abschnitt~\ref{anbindbarkeit_an_externe_software} durchgeführte schnelle Umsetzung eines Adapters 
zur Anbindung der \emph{MANVSuite} an eine Weboberfläche konnte gezeigt werden, dass durch die hohe Modularität
und Einsatz der Komponententechnologie \emph{Corba} ein hohes Maß an Flexibilität und Adaptierbarkeit der
Gesamtlösung erreicht wird und eine Anbindung an bestehende Systeme mit geringem Aufwand realisiert werden kann.
Eine genauere Analyse und Evaluation der \emph{MANVSuite} inkl. der beschriebenen Schnittstellen findet sich
in \cite{Jan}. In dieser verwandten Arbeit wird unter anderem gezeigt, wie eine Anbindung der \emph{MANVSuite}
an ein Mobiltelefon des Herstellers Nokia erfolgen kann.\\
\\
Als besonders positiv ist die Skalierbarkeit der Gesamtlösung zu bewerten: Die Verbindung der einzelnen Komponenten 
untereinander ist netzwerkfähig, im Extremfall kann jede
Komponente auf einem eigenen Rechner betrieben werden. Hierbei skaliert die Anforderung an die Hardware mit
der Größe des zu verwaltenden Netzwerkes. Im Prinzip ist jedes linuxfähige Gerät mit mindestens 16~MB Speicher
und einer \emph{USB-Schnittstelle}, für die eine Implementierung der \emph{Java Virtual Machine}\footnote{Hierbei muss
es sich nicht zwangsläufig um die Implementierung von Oracle/Sun handeln. Die \emph{MANVSuite} kann z.B. auch
mit der \emph{JVM} des \emph{Eclipse-Projektes} ausgeführt werden.} verfügbar ist, in der Lage, als \emph{Server}
für das Sensornetzwerk zu dienen. Dies kann z.B. ein \emph{Android} basiertes Handy oder ein offener
\emph{DSL}-Router in Verbindung mit dem Betriebssystem \emph{OpenWRT} sein. Sollte keine \emph{JVM} für die
Plattform verfügbar sein, besteht sogar die Möglichkeit, dass nur der in \emph{Python} geschriebene,
hardwarespezifische Teil des \emph{MANVConnectors} auf dem Gerät selbst ausgeführt wird\footnote{Die Implementierung
in \emph{Python} erfolgte lediglich aus Gründen der Einfachheit. Sollte kein \emph{Python} verfügbar sein,
kann dieser Teil in jeder anderen Programmiersprache, die \emph{Socketverbindungen} erlaubt, neu implementiert werden.}. 
Alle restlichen Teile der \emph{MANVSuite} können nachgelagert auf einem Server im Internet ausgeführt werden; 
die Verbindung der beiden Teile erfolgt über eine einfache \emph{Socketverbindung}. Diese Lösung bietet sich insbesondere 
für die Überwachung einzelner Patienten als eine Art intelligenter Hausnotruf an.\\
\\
Vom bisherigen Stand der Technik hebt sich die hier vorgestellte Lösung in vielen Punkten ab:

\nomenclature{Open Source}{Softwareprodukt, bei dem der Quellcode frei verfügbar ist.}
\nomenclature{GPIO}{General Purpose Input Output}
\begin{itemize}
    \item{Konsequenter Einsatz offener Standards:} In der gesamten Arbeit wurde konsequent auf den
          Einsatz offener Standards gesetzt. Bei der Verwendung von externen Softwarepaketen wurde
          darauf geachtet, dass es sich um \emph{Open-Source} Produkte handelt. Zur Implementierung
          wurde keinerlei proprietäre Technologie eingesetzt.

    \item{Kapselung aller plattformspezifischen Designentscheidungen:} Alle Plattformabhängigen
          Designentscheidungen wie die Ansteuerung der \emph{ZigBee}-Module oder der
          physische Zugriff auf die serielle Schnittstelle sind bewusst austauschbar gehalten.
          Hierzu wurden diese in Komponenten gekapselt und die Schnittstelle zwischen diesen
          Komponenten wurde standardisiert. Bei Änderung einer dieser Entscheidungen (z.B.
          weil ein anderes Funkprotokoll oder ein anderes Betriebssystem zum Einsatz kommen soll)
          können diese einfach ausgetauscht werden, ohne den Rest des Systems anpassen zu 
          müssen.

    \item{Geringe Stückkosten:} Bei der Auswahl des Funkprotokolls wurde insbesondere auch auf den
          Preis der verfügbaren Hardwaremodule geachtet. Dieser bewegt sich selbst bei der
          Einzelabnahme der Module in einem Bereich unter 15\EUR{}. Auf den Einsatz teurer
          Basisbauteile wie z.B. \emph{TinyOS}-basierten Sensoren wurde konsequent verzichtet.

    \item{Geringe Leistungsaufnahme ("`Stromverbrauch"'):} Durch den Einsatz des von \emph{ZigBee}
          angebotenen Energiesparmodus konnte der mittlere Stromverbrauch auf einen Wert von ca. 
          1,86~mA gesenkt werden.

    \item{Einfache Implementierbarkeit:} Zur Anbindung des \emph{ZigBit}-Moduls an den \emph{Erste-Hilfe-Sensor}
          sind lediglich vier Signalleitungen notwendig: \emph{UART\_RX}, \emph{UART\_TX}, 
          \emph{UART\_CTS}\footnote{UART\_RTS sollte auf Seite des \emph{ZigBit}-Moduls auf \emph{GND}
              geführt werden.} sowie eine \emph{GPIO}-Leitung, die mit dem \emph{RESET}-Eingan des
          \emph{ZigBit}-Moduls verbunden ist (Details zur Implementierung in die Platine können dem
          Schaltplan in Anhang~\ref{anhang_diagramme} entnommen werden). 

    \item{Plattformunabhängigkeit:} Durch den Einsatz plattformunabhängiger Programmiersprachen wie \emph{Java}
          und \emph{Python} sowie durch Hardwarevirtualisierung ist die vorgestellte Lösung auf
          einer breiten Palette von Systemen einsatzfähig.

    \item{Kompaktheit:} Die hier vorgestellte Prototyp-Platine besitzt eine Abmessung von nur
          50~mm auf 56~mm und bietet noch viel Potential zur weiteren Optimierung der Größe (z.B. durch
          die beidseitige Platzierung von Bauteilen oder Fertigung in 4-Schicht-Bauweise).
\end{itemize}

In dieser Arbeit nicht betrachtet wurde das Thema Datensicherheit und Verschlüsselung, welches alleine bereits
umfangreich genug wäre, eine eigene Diplomarbeit zu rechtfertigen. Ursprünglich war geplant, die hardwareseitige
Verschlüsselung der \emph{ZigBit}-Module einzusetzen. Leider ist diese Funktion jedoch nicht in die 
\emph{SerialNet}-Firmware integriert. Daher wäre für die Verwendung der Verschlüsselung eine Neuimplementierung der
\emph{ZigBit-Firmware} notwendig, was den Rahmen dieser Arbeit sprengen würde. Daher beschränkt 
sich diese Arbeit darauf, in Abschnitt~\ref{Sicherheit} eine kurze Analyse des Problems und einige
Vorschläge zur Lösung vorzustellen. Diese können dann bei einer späteren Umsetzung in ein Produkt als Grundlage
dafür dienen, die notwendigen Sicherheitsmaßnahmen zu implementieren.

%%%------------------------------------------------ZUSAMMENFASSUNG UND AUSBLICK------------------------------------------------------

\chapter{Zusammenfassung und Ausblick}
***Beispiel wie ein Zitat auf Armins~\citep{Bolz} Buch funktioniert***
\section{Verbesserung der Sicherheit}
-> Eigene Firmware für ZigBit

\section{Verbesserung der Reichweite und Störsicherheit}
\subsection{Anderer Frequenzbereich}
\subsection{Richtantennen für Router}
\subsection{Adaptive Pfadkosten für Meshrouting}


%%%------------------------------------------------ANHANG---------------------------------------------------------------

\begin{appendix}

%%%------------------------------------------------STAND DER TECHNIK------------------------------------------------------------------
\chapter{Pläne und Diagramme}\label{anhang_diagramme}
    \Abbildungps{!htb}{0.4}{diagramme/schaltplan_usbstick.pdf}{label}{Schaltplan des \emph{USB-Sticks}.}
    \Abbildungps{!htb}{1.05}{diagramme/schaltplan_manvnode.pdf}{label}{Schaltplan der \emph{MANVNode}.}
    \Abbildungps{!htb}{0.75}{diagramme/platine_manvnode_top.pdf}{label}{Oberseite der Platine der MANVNode.}
    \Abbildungps{!htb}{0.75}{diagramme/platine_manvnode_bottom.pdf}{label}{Unterseite der Platine der MANVNode.}
    \begin{sidewaysfigure}
            \includegraphics*[width=1\textwidth]{uml_diagramme/manvconnector_uml.pdf}
            \caption{Klassendiagramm des MANV-Connectors.}
            \label{klassendiagramm_manvconnector}
    \end{sidewaysfigure}

\chapter{Beschreibung der beigelegten Software}\label{anhang_beschreibung_software}
    \section{MANV-Connector}

    Im Verzeichnis "`connector"' auf der beigelegten CD befinden sich
    Quellcode und Binaries des \emph{MANV-Connectors}. Zum Kompilieren
    und Starten des Connectors wird das Script "`start.sh"' verwendet.

    \subsection{Starten des Connectors}

    Zum Starten des Connectors muss das Skript "`start.sh"' mit dem Parameter
    "`start"' aufgerufen werden. Evtl. müssen in diesem Skript noch Pfade
    und die IP-Adresse des \emph{MANV-Servers} angepasst werden. 
    Dieser Server muss vor dem \emph{MANV-Connector} gestartet werden.

    \subsection{Kompilieren des Connectors}

    Zum Kompilieren der beigelegten Software wird ebenfalls das "`start.sh"'
    Skript verwendet. Hierzu wird dieses mit dem Parameter "`build"' aufgerufen.
    Zum erfolgreichen Kompilieren sind die \emph{IDL}-Dateien der \emph{MANV-Suite}
    von Jan Tepelmann erforderlich, die auf der CD im Verzeichnis "`jan"'
    mitgeliefert werden. Evtl. müssen hierfür die Pfade im "`start.sh"'-Skript 
    angepasst werden.

    \subsection{Verzeichnisstruktur des Quellcodes}

    Die Verzeichnisstruktur des Quellcodes folgt dem Standard für \emph{Java-Packages}.
    Der Quellcode ist auf mehrere Dateien und Unververzeichnisse aufgeteilt
    und findet sich im Verzeichnis \emph{connector/edu/kit/ibt/manv/connector}.
    
    \begin{table}
        \begin{tabular}{ll}
            \textbf{Verzeichnis} & \textbf{Beschreibung}\\

           commands/ & Befehle, die an den \emph{USB-Stick} gesendet werden.\\
           corba/    & Implementierung der \emph{CORBA}-Schnittstelle.\\
           events/   & Ereignisse, die von dem \emph{USB-Stick} empfangen wurden.\\
           results/  & Antworten auf gesendete Befehle.\\
           lib/      & Klassen, die thematisch nicht in andere Pakete passen.\\
        \end{tabular}
        \caption{Struktur des Quellcodeverzeichnisses}
    \end{table}

    \subsubsection{Inhalt des Verzeichnises \emph{commands/}}

    Innerhalb des \emph{commands/}-Verzeichnisses findet sich die Datei \emph{MANVCommands.java},
    in der die Kasse \emph{MANVCommand} spezifiziert wird. Objekte dieser Klasse kapseln jeweils
    einen \emph{SerialNet}-AT-Befehl (vergleiche Tabelle~\ref{serialnet_befehle}), der an den
    \emph{USB-Stick} gesendet wird. Dieser Befehl wird als einfacher ASCII-String an den Konstruktor
    übergeben. Optional kann bei wichtigen Befehlen noch eine Priorität mitgegeben werden;
    je höher diese Priorität ist, desto eher wird dieser Befehl an den \emph{USB-Stick}
    gesendet. Wenn der Befehl abgearbeitet wurde, wird in dem Objekt das Ergebniss
    der Ausführung gespeichert, welches mit der Methode \emph{MANVResult.getResult()}
    abgerufen werden kann.

    \subsubsection{Inhalt des Verzeichnises \emph{corba/}}

    In dem Verzeichniss \emph{corba/} findet sich der Glue-Code der Implementierung der \emph{Corba}-Schnittstelle
    zum \emph{MANV}-Server. Die einzelnen Dateien haben folgende Funktion:

    \begin{itemize}
        \item{CommandsImpl.java:} Definiert, wie auf vom Server empfangene Befehle reagiert werden soll.
                                  Leitet empfangene Befehle an das \emph{ZigBee}-Netzwerk weiter.
        \item{MessageContainer.java:} Container zur Auslieferung von \emph{ZigBee}-Nachrichten in \emph{Corba}-Nachrichten.
                                      Für jeden Nachrichtentyp gibt es eine spezielle Unterklasse 
                                      (\emph{CorbaDataMessageContainer},
                                      \emph{CorbaEventMessageContainer}). Der Inhalt der Container wird von der Methode 
                                      \emph{createCorbaMessages()} der Unterklassen von \emph{events/MANVEvent.java}
                                      konstruiert.
        \item{CobaSender.java:} Thread, der empfangene \emph{ZigBee}-Ereignisse an den \emph{MANV-Server} ausliefert.
                                Diese Ereignisse werden aus einer Queue (\emph{eventQueue}) entnommen und sind
                                bereits in \emph{MessageContainer} verpackt.
    \end{itemize}

    \subsubsection{Inhalt des Verzeichnisses \emph{events/}}
    
    Im Verzeichnis \emph{events/} befinden sich die Klassen, welche Ereignisse repräsentieren, die von dem 
    \emph{ZigBee}-Netzwerk empfangen wurden. Alle Ereignis-Klassen sind von der Klasse \emph{MANVEvent} abgeleitet 
    und besitzen damit automatisch eine Priorität. Über diese Priorität kann eingestellt werden, in welcher Reihenfolge 
    empfangene Ereignisse abgearbeitet werden sollen. So ist es z.B. sinnvoll, dass Ereignisse, welche einen Alarm 
    darstellen, eine höhere Priorität haben als eine Mitteilung, dass ein neuer Knoten im Netzwerk vorhanden ist. 
    Nachfolgend eine Übersicht über die Funktion der Methoden der Klasse \emph{MANVEvent}:

    \begin{itemize}
            \item{\emph{isResult()}:} Gibt zurück, ob es sich bei dem Ereignis um das Ergebnis eines gesendeten Befehls handelt. 
            \item{\emph{isImportant()}:} Unterscheidet wichtige von unwichtigen Ereignissen. Nur wichtige Ereignisse werden über 
                                  die \emph{CORBA}-Schnittstelle an den \emph{MANV-Server} weitergeleitet. Unwichtige 
                                  Ereignisse sind nur für die interne Verarbeitung im \emph{MANV-Connector} relevant. 
            \item{\emph{fromString()}:}  Factory-Methode, die aus einem \emph{SerialNet}-String ein \emph{MANV\-Event} erstellt. 
                                  Liefer automatisch eine Instanz einer Unterklasse von \emph{MANV\-Event} zurück.
            \item{\emph{createCorbaMessages()}:} Bereitet ein \emph{MANVEvent} für den Transport über die 
                                          \emph{CORBA}-Schnittstelle zum \emph{MANV-Server} vor. Hierzu werden ein oder 
                                          mehrere \emph{CorbaEvent}-Objekte erstellt, die wiederum in einen
                                          passenden \emph{CorbaContainer} verpackt werden.
            \item{\emph{getRaw()}:} Liefert den rohen \emph{SerialNet}-String zurück, auf Basis dessen das \emph{MANVEvent}-Objekt 
                             erzeugt wurde; dies ist vor allem für Debugging-Zwecke interessant.
    \end{itemize}

    Von der Klasse \emph{MANVEvent} sind folgende Unterklassen abgeleitet:

    \begin{itemize}
        \item{\emph{MANVChildJoined}:} Eine neue \emph{MANVNode} ist dem Netzwerk beigetreten.
        \item{\emph{MANVChildLost}:} Eine \emph{MANVNode} hat die Verbindung zum Netzwerk verloren.
        \item{\emph{MANVStatusMessage}:} Eine Datennachricht, die Statusinformationen (Alarmzustand, Puls, Atmung) enthält,
                                  wurde empfangen.
        \item{\emph{MANVDataReceived}:} Eine Datennachricht, die keine Statusinformationen enthält, wurde empfangen.
    \end{itemize}

    \subsubsection{Inhalt der Verzeichnises \emph{results/}}

    In dem Verzeichnis \emph{results/} befinden sich alle Klassen, die auf gesendete \emph{SerialNet}-Befehle empfangene 
    Antworten repräsentieren. Diese sind von der Basisklasse \emph{MANVResult} abgeleitet, welche wiederum eine Unterklasse 
    von \emph{MANVEvent} ist.  Damit besitzt diese Klasse alle Eigenschaften und Methoden, die auch die 
    Klasse \emph{MANVEvent} besitzt; insbesondere handelt es sich hierbei selbst wieder um eine priorisierte Nachricht. 
    Zusätzlich werden folgende Methoden neu eingeführt:

    \begin{itemize}
        \item{\emph{getData()}:} Liefert die von der \emph{SerialNet-Firmware} empfangenen Daten in Stringform zurück. 
        \item{\emph{isComposite()}:} Gibt zurück, ob es sich um eine komplexe Antwort (vgl. 
                            Abschnitt~\ref{analyse_kommunikation_mit_zigbit_modul}) handelt. In diesem Fall befindet 
                            sich in dem Attribut \emph{subResult} eine gültige Referenz auf ein Unterergebnis.
        \item{\emph{getChildList()}:} Falls mit dem gesendeten \emph{SerialNet}-Befehl nach den assoziierten Kindknoten gefragt wurde,
                             gibt diese Methode eine Liste dieser Kindknoten zurück.
    \end{itemize}

    Neben den beiden Statusinformationen \emph{OK} und \emph{ERROR} kann die \emph{SerialNet-Firmware} auch 
    kompliziertere Antworten zurückliefern. Diese werden von folgenden Unterklassen von \emph{MANVEvent} repräsentiert:
        
    \begin{itemize}
        \item{\emph{MANVChildrenList}:} Repräsentiert eine Liste von Kindknoten, die von der \emph{SerialNet-Firmware}
                                 empfangen wurde.
        \item{\emph{MANVGsn}:} Liefert die \emph{Global Subscriber Number} (GSN) einer Node zurück. Entspricht in etwa einer 
                        \emph{MAC}-Adresse.
    \end{itemize}

    \subsubsection{Inhalt des Verzeichnisses \emph{lib/}}

    In dem Verzeichnis \emph{lib/} finden sich alle die Klassen, welche nicht in obiges Einordnungsschema passen. 
    Hierbei handelt es sich insbesondere um Helferklassen, die in den obigen Klassen verwendet werden:

    \begin{itemize}
        \item{\emph{iZigBit}:} Interface für einen entfernten \emph{ZigBit}-Knoten. Diese werden nochmal weiter
                        in vollwertige Knoten (\emph{ZigBit}, im wesentlichen also die \emph{MANV-Nodes})
                        und Knoten, von denen nur Daten empfangen werden können (\emph{readonlyZigBit)})
                        unterschieden.
        \item{\emph{MANVPrioritized}:} Allgemeine Oberklasse für priorisierte Entitäten wie z.B. \emph{MANVCommand},
                                \emph{MANVEvent} oder \emph{MANVResult}. Diese Entitäten enthalten 
                                sowohl eine numerische Priorität als auch eine \emph{compareTo()}-Operation,
                                über die eine Sortierung vorgenommen werden kann.
        \item{\emph{SocketReader}:} Definiert einen Thread, der von dem Socket liest, der die Verbindung zum 
                                    \emph{MANV-USB-Stick} beinhaltet. Dieser Thread empfängt Ereignisse und 
                                    Ergebnisse, wandelt sie in Instanzen vom Typ \emph{MANVEvent} und 
                                    \emph{MANVResult} um und sortiert
                             in die Ereigniswarteschlange (\emph{eventQueue}) ein. Empfangene Ergebnisse werden
                             automatisch dem passenden Objekt vom Typ \emph{MANVCommand} zugeordnet.
        \item{\emph{SocketWriter}:} Definiert einen Thread, der zu sendende Befehle aus der Befehlswarteschlange
                             (\emph{commandQueue}) entnimmt und über die Socketverbindung zum \emph{MANV-USB-Stick}
                             an das \emph{ZigBee}-Netzwerk weiterleitet. Sobald ein Befehl abgearbeitet ist,
                             wird er von dem \emph{SocketWriter}-Thread automatisch in die Ergebniswarteschlange 
                             (\emph{resultQueue}) einsortiert.
    \end{itemize}


    \paragraph{Interface eines ZigBit-Knotens}

    Das Interface \emph{iZigBit} definiert folgende Methoden:

    \begin{itemize}
        \item{\emph{getNodeID()}:} Liefert die logische Adresse (\emph{WSRC}) eines Knotens zurück.
        \item{\emph{getMacID()}:} Liefert die Hardware-Adresse (\emph{GSN}) eines Knotens zurück.
        \item{\emph{sendData()}:} Sendet Daten an einen Knoten. Blockiert so lange, bis die Bearbeitung des
                           Befehls abgeschlossen ist.
        \item{\emph{isendData()}:} Senden von Daten an einen Knoten. Kehrt sofort zurück, ohne auf die Bearbeitung
                            des Befehls zu warten.
        \item{\emph{toggleAlertStatus()}:} Umschalten des Alarmstatus einer \emph{MANV-Node}.
        \item{\emph{disableAlert()}:} Deaktivieren des Alarms einer \emph{MANV-Node}.
        \item{\emph{enableAlert()}:} Aktivieren des Alarms einer \emph{MANV-Node}.
        \item{\emph{muteAlert()}:} Alarm stummschalten.
        \item{\emph{itoggleAlertStatus()}, \emph{idisableAlertStatus()}, \emph{ienableAlert()}, \emph{imuteAlert()}}: Selbe Funktionalität
            wie die Versionen ohne i, jedoch ohne Warten auf Abarbeitung (analog zu \emph{sendData()} / \emph{isendData()}).
    \end{itemize}

    \section{Serial-To-Socket}

    In dem Verzeichnis \emph{serial\_to\_socket} befinden sich eine Reihe von Python-Skripten, die die
    Verbindung zwischen \emph{MANVConnector} und \emph{MANV-USB-Stick} herstellen. Hierzu 
    wird zunächst mit Hilfe des Skriptes \emph{serial\_to\_socket.py} eine automatische Suche nach dem 
    \emph{MANV-USB-Sticks} (vgl. Abschnitt~\ref{usb_stick_erkennen}) durchgeführt. Der Status
    dieser Suche wird graphisch ausgegeben. Sobald ein \emph{MANV-USB-Stick} erfolgreich
    erkannt wurde, wird das Skript \emph{serial\_to\_socket/start.sh} gestartet und die
    Adresse des seriellen Ports als Parameter übergeben. In diesem Skript wird
    dann unter anderem das Pythonskript "`serial\_to\_socket.py"' gestart. Dieses leitet einen
    angegebenen seriellen Port auf den Netzwerkport \emph{4711} um, über den dann der
    \emph{MANVConnector} eine Verbindung zu dem \emph{USB-Stick} herstellen kann.
    
    \section{Firmware}

    In dem Verzeichnis \emph{firmware/} befindet sich der Quellcode der \emph{Firmware} für die \emph{MANVNode}.
    Diese besteht aus folgenden Dateien:

    \begin{itemize}
        \item{\emph{main.c}:} Einstiegspunkt der \emph{Firmware}. Initialisiert alle Ringpuffer, \emph{Interrupts} und 
                       die \emph{Peripherie} des \emph{ADuC}-Mikrocontrollers. Ruft danach die 
                       Hauptscheife des \emph{ZigBit}-Treibers auf.
        \item{\emph{isr.c}:}  Definition der verwendeten \emph{Interrupt-Service-Routinen}, also Behandlungsroutinen
                       für auftretende \emph{Interrupts}. Ist unter anderem für die Ansteuerung des \emph{UART} 
                       verantwortlich.
        \item{\emph{random.c}:} Implementierung eines sehr einfachen Pseudo-Zufallszahlengenerators.
        \item{\emph{ringpuffer.c}:} Implementierung der Datenstruktur "`Ringpuffer"'. Diese kann 5 Zeilen
                             von jeweils 80 Zeichen Länge speichern. Werden mehr als 5 Zeilen in die
                             Datenstruktur gelegt, dann wird das älteste Element überschrieben.
                             Ringpuffer werden unter anderem zur Implementierung des \emph{UART}-Treibers
                             sowie zur Zwischenspeicherung von empfangenen Befehlen zur späteren
                             Abarbeitung verwendet.
        \item{\emph{uart.c}:} Beinhaltet die Routinen, die zur Interaktion mit dem als \emph{ISR} implementierten
                       Treiber für den \emph{UART} notwendig sind. Im wesentlichen handelt es sich um 
                       Routinen, die die entsprechenden Ringpuffer befüllen oder auslesen.
        \item{\emph{zigbit.c}:} Beispiel einer Interaktion mit dem \emph{ZigBit}-Modul. Enthält unter anderem Routinen
                         zur Befehlsverareitung, Extraktion von Rückgabewerten, Initialisierung des \emph{ZigBit}-Moduls
                         sowie ein Beispiel, wie die Übertragung von Statusinformationen realisiert werden kann.
    \end{itemize}

    Bei der \emph{Firmware} handelt es sich um eine Demonstration, wie die Kommunikation zwischen \emph{ADuC}-Mikrocontroller
    und \emph{ZigBit}-Modul realisiert werden kann. Zur späteren Verwendung auf dem \emph{Erste-Hilfe-Sensor} müssen 
    einige Anpassungen vorgenommen werden. Hierzu wird folgende Vorgehensweise vorgeschlagen:

    \begin{enumerate}
        \item Übernahme der \emph{UART}-spezifischen Routinen aus \emph{isr.c} und \emph{uart.c}. 
        \item Übernahme der Datenstruktur \emph{Ringpuffer} aus \emph{ringpuffer.c}.
        \item Übernahme der Funktionen \emph{extractStatus()}, \emph{waitForStatus()}, \emph{isCommand()}
              und \emph{parseCommand()} aus der Datei \emph{zigbit.c}.
        \item Anpassung der Funktion \emph{initZigBit()} aus der Datei \emph{zigbit.c}:
              Zu beachten ist hierbei Zeile 154: Hier wird eine \emph{GPIO}-Leitung des \emph{ADuC} Mikrocontrollers 
              dazu verwendet,
              die Resetleitung des \emph{ZigBit}-Moduls zu betätigen. Dies ist für die Synchronisierung
              zwingend erforderlich und muss im Hardwarelayout entsprechend berücksichtigt werden.
              Nach Zeile 154 sollten alle Zeilen mit Zugriff auf \emph{GP2DAT} entfernt werden, 
              da diese dazu dienen, die \emph{LEDs} der \emph{MANVNode} anzusteuern. Diese \emph{LEDs} sind
              auf dem Erste-Hilfe-Sensor nicht vorhanden. Die Aufrufe von \emph{sleep()} sind zur Synchronisierung
              zwingend erforderlich und dürfen nicht entfernt werden.
        \item Anpassung der Funktion \emph{zigBitLoop()} aus der Datei \emph{zigbit.c} Hier werden aktuell noch 
              Dummy-Informationen übertragen, die durch einen Zufallszahlengenerator erzeugt werden. Diese
              müssen durch die tatsächlich gemessenen Werte des Erste-Hilfe-Sensors ersetzt werden.
    \end{enumerate}                


    \section{MANV-Web}
    Im Verzeichnis \emph{manvweb/} befinden sich die Quelldateien des Webinterfaces der
    \emph{MANVSuite}. Dabei verteilt sich die Funktionalität auf folgende Unterdateien:

    \begin{itemize}
        \item{\emph{MANVWeb}:} Einstiegsklasse des Webservers. Stellt die \emph{CORBA}-Verbindung her
                        und übergibt dann die Steuerung an die Klasse \emph{MANVWebserver}.
        \item{\emph{MANVWebserver}:} Implementierung des Webervers. Öffnet einen Serversocket und
                              wartet auf eingehende Verbindungen. Für jede eingehende Verbindung
                              wird nun ein neuer Thread gestartet. In diesem Thread
                              wird die Eingabe geparsed und eine entsprechende Reaktion
                              aufgerufen. 
        \item{\emph{ToQueries}:} Adapterklasse zum \emph{CORBA}-Zugriff auf den \emph{MANV-Server}. 
    \end{itemize}

    \section{Platine}

    In dem Verzeichnis \emph{platine/} befinden sich die \emph{Eagle}-Dateien, in 
    denen Schaltplan und Platinenlayout des \emph{MANV-USB-Sticks} und der
    \emph{MANV-Node} spezifiziert sind.

    \section{Simulator}

    In dem Verzeichnis \emph{simulator/} befindet sich ein in Python programmierter
    \emph{MANVNode}-Simulator. Dieser kann dazu verwendet werden, mit Hilfe
    eines \emph{MANV-USB-Sticks}, \emph{MANVNodes} zu simulieren. Beim Start des
    Simulators muss diesem die Adresse des seriellen Ports, auf dem sich der
    \emph{USB-Stick} befindet, übergeben werden.

    \section{Zigbit Firmware}
    In dem Verzeichnis \emph{ZigBit-Firmware} befindet sich eine aktuelle Kopie der
    \emph{SerialNet-Firmware} inkl. dem zum Programmieren der \emph{ZigBit}-Module
    nötigen Bootloader in einer Version für Windows.

\end{appendix}



%%%------------------------------------------------KOMPLETTES LITERATURVERZEICHNIS NACH NORM------------------------------------------
\nocite{*} % Alle Einträge anzeigen
\bibliography{Diplomarbeit}
\bibliographystyle{itmalpha}
\addcontentsline{toc}{chapter}{Literaturverzeichnis}

\end{document}

